\subsection{相对于模的范数与迹}

定义1. 设 M 是一个 A-模，且作为 K-模具有有限基。对所有 \(a \in \mathbf{A}\) ，设 \(a_{\mathbf{M}}\) 为 K-模 M 的自同态 \(x \mapsto ax\) 。 \(a_{\mathbf{M}}\) 的迹、行列式和特征多项式分别称为 \(z\) 相对于 M 的迹、范数和特征多项式。

因此，a 的迹和范数是 K 的元素，分别记为 \(\mathrm{Tr}_{\mathbf{M}/\mathbf{K}}(a)\) 和 \(\mathrm{N}_{\mathbf{M}/\mathbf{K}}(a)\) ；a 的特征多项式是 K{[}X{]} 的元素，记作 \(\mathrm{Pc}_{\mathbf{M}/\mathbf{K}}(a;\mathbf{X})\) 。当没有混淆风险时，我们在上述记号中省略 K。

由自同态的迹和行列式的性质（II，第4节，第3小节及第8节，第1小节），我们得到关系式

\[
\operatorname{Tr} _ {\mathbf {M}} (a + a ^ {\prime}) = \operatorname{Tr} _ {\mathbf {M}} (a) + \operatorname{Tr} _ {\mathbf {M}} (a ^ {\prime})\tag{1}
\]

\begin{enumerate}
\def\labelenumi{(\arabic{enumi})}
\setcounter{enumi}{1}
\tightlist
\item
\end{enumerate}

\[
\operatorname{Tr} _ {\mathbf {M}} (a a ^ {\prime}) = \operatorname{Tr} _ {\mathbf {M}} (a ^ {\prime} a)\tag{3}
\]

\[
\mathrm{N} _ {\mathbf {M}} (a a ^ {\prime}) = \mathrm{N} _ {\mathbf {M}} (a) \mathrm{N} _ {\mathbf {M}} (a ^ {\prime})
\]

对 A 中所有 \(a, a'\) 成立。

设 \((e_i)_{1\leq i\leq n}\) 是 K-模 M 的一组基，且 \((m_{ij}(a))\) 是自同态 \(a_{\mathbf{M}}\) 关于这组基的矩阵。函数 \(m_{ij}\) 是K-模A上的线性形式，且

\begin{enumerate}
\def\labelenumi{(\arabic{enumi})}
\setcounter{enumi}{3}
\tightlist
\item
\end{enumerate}

\[
\operatorname{Tr} _ {\mathbf {M}} (a) = \sum_ {i = 1} ^ {n} m _ {i i} (a)\tag{5}
\]

\[
\mathrm{N} _ {\mathbf {M}} (a) = \det (m _ {i j} (a))\tag{6}
\]

\[
\operatorname{Pc} (a; \mathbf {X}) = \det (\delta_ {i j} \mathbf {X} - m _ {i j} (a)).
\]

由计算行列式的方法（第8节，第11小节，公式(50)）可知

\[
\mathrm{Pc} _ {\mathbf {M}} (a; \mathbf {X}) = \mathbf {X} ^ {n} + c _ {1} \mathbf {X} ^ {n - 1} + \dots + c _ {n}\tag{7}
\]

其中

\[
c _ {1} = - \operatorname{Tr} _ {\mathbf {M}} (a), \quad c _ {n} = (- 1) ^ {n} \mathrm{N} _ {\mathbf {M}} (a).\tag{8}
\]

对于 \(\lambda \in \mathbf{K}\) ，

\[
\operatorname{Tr} _ {\mathbf {M}} (\lambda) = n. \lambda , \quad \mathrm{N} _ {\mathbf {M}} (\lambda) = \lambda^ {n}, \quad \operatorname{Pc} _ {\mathbf {M}} (\lambda ; \mathbf {X}) = (\mathbf {X} - \lambda) ^ {n}.\tag{9}
\]

设 \(\mathbf{K}'\) 是一个交换 K-代数。设 \(\mathbf{M}' = \mathbf{K}' \otimes_{\mathbf{K}} \mathbf{M}\) 和 \(\mathbf{A}' = \mathbf{K}' \otimes_{\mathbf{K}} \mathbf{A}\) ，于是 \(\mathbf{M}'\) 具有一个 A'-模结构（§4，例 2）。作为 \(\mathbf{K}'\) -模， \(\mathbf{M}'\) 具有由 \(1 \otimes e_i\) ( \(1 \leqslant i \leqslant n\) ) 组成的基，且 \(a_{\mathbf{M}}\) 关于 \((e_i)\) 的矩阵等于 \((1 \otimes a)_{\mathbf{M}'}\) 关于 \((e_i')\) 的矩阵。则

\[
\operatorname{Tr} _ {\mathbf {M} ^ {\prime}} (1 \otimes a) = \operatorname{Tr} _ {\mathbf {M}} (a). 1, \quad \mathrm{N} _ {\mathbf {M} ^ {\prime}} (1 \otimes a) = \mathrm{N} _ {\mathbf {M}} (a). 1\tag{12}
\]

\[
\mathrm{Pc} _ {\mathbf {M} ^ {\prime}} (1 \otimes a; \mathbf {X}) = \mathrm{Pc} _ {\mathbf {M}} (a; \mathbf {X}). 1
\]

对所有 \(a \in \mathbf{A}\) 成立，其中 1 表示 \(\mathbf{K}'\) 的单位元。特别地，若我们取 \(\mathbf{K}' = \mathbf{K}[\mathbf{X}]\) ，则

\[
\mathrm{Pc} _ {\mathbf {M} / \mathbf {K}} (a; \mathbf {X}) = \mathrm{N} _ {\mathbf {M} [ \mathbf {X} ] / \mathbf {K} [ \mathbf {X} ]} (\mathbf {X} - a).\tag{13}
\]

\subsection{相对于模的范数与迹的性质}

若 \(\mathbf{M}\) 和 \(\mathbf{M}'\) 是两个同构的 A-模，且在 \(\mathbf{K}\) 上具有有限基，那么对所有 \(a \in \mathbf{A}\) ，

\[
\operatorname{Tr} _ {\mathbf {M} ^ {\prime}} (a) = \operatorname{Tr} _ {\mathbf {M}} (a), \quad \mathrm{N} _ {\mathbf {M} ^ {\prime}} (a) = \mathrm{N} _ {\mathbf {M}} (a), \quad \operatorname{Pc} _ {\mathbf {M} ^ {\prime}} (a; \mathbf {X}) = \operatorname{Pc} _ {\mathbf {M}} (a; \mathbf {X})\tag{14}
\]

因为若 \(f\) 是从 \(\mathbf{M}\) 到 \(\mathbf{M}'\) 的同构，则 \(a_{\mathbf{M}}\) 关于基 \(\mathbf{B}\) （ \(\mathbf{M}\) 在 \(\mathbf{K}\) 上的一个基）的矩阵，与 \(a_{\mathbf{M}'}\) 关于基 \(f(\mathbf{B})\) （ \(\mathbf{M}'\) 的基）的矩阵相同。

命题 1. 设 \(\mathbf{M} = \mathbf{M}_0 \supset \mathbf{M}_1 \supset \cdots \supset \mathbf{M}_r = \{0\}\) 是 A-模 \(\mathbf{M}\) 的子模的一个递减序列，使得每个 K-模 \(\mathbf{P}_i = \mathbf{M}_{i-1} / \mathbf{M}_i\) ( \(1 \leqslant i \leqslant r\) ) 都具有有限基。则 K-模 \(\mathbf{M}\) 具有有限基，并且

\[
\operatorname{Tr} _ {\mathbf {M}} (a) = \sum_ {i = 1} ^ {r} \operatorname{Tr} _ {\mathrm{P} _ {i}} (a), \quad \mathrm{N} _ {\mathbf {M}} (a) = \prod_ {i = 1} ^ {r} \mathrm{N} _ {\mathrm{P} _ {i}} (a)\tag{15}
\]

\[
\mathrm{Pc} _ {\mathbf {M}} (a; \mathbf {X}) = \prod_ {i = 1} ^ {r} \mathrm{Pc} _ {\mathbf {P} _ {i}} (a; \mathbf {X}).
\]

设 \(B_{i}^{\prime}\) 是 \(P_{i}\) 在 K 上的一组基；则 \(B_{i}\) 作为 \(B_{i}^{\prime}\) (模 \(M_{i}\) ) 的一个代表系，是 K-模 \(M_{i}\) 在 K-模 \(M_{i-1}\) 中的一个补子模的基（II，§1，第11号，命题21）。诸 \(B_{i}\) ( \(1 \leqslant i \leqslant r\) ) 的并集 B 是 M 在 K 上的一组基。设 \(X_{ii}\) 是自同态 \(a_{P_{i}}\) 关于基 \(B_{i}^{\prime}\) 的矩阵。立即可以看出， \(a_{M}\) 关于 B 的矩阵具有如下形式

\[
\left( \begin{array}{c c c c c} X _ {r r} & X _ {r, r - 1} & \ldots & X _ {r, 1} \\ 0 & X _ {r - 1, r - 1} & \ldots & X _ {r - 1, 1} \\ \cdot & \cdot & \cdot & \cdot & \cdot \\ 0 & 0 & \ldots & X _ {1 1} \end{array} \right)
\]

于是命题由第 1 号中的公式 (4)、(5)、(6) 以及 § 8 第 6 号中的公式 (31) 得出。

命题 2. 设 A、A' 为两个 K-代数，M 为 A-模，M' 为 A'-模。假设 M 与 M' 是维数分别为 n 与 n' 的自由 K-模，并将 M ⊗K M' 视为 (A ⊗K A')-模（§ 4，第 3 号，例 2）。那么，对于 a \ensuremath{\in} A 与 a' \ensuremath{\in} A'，

\begin{enumerate}
\def\labelenumi{(\arabic{enumi})}
\setcounter{enumi}{15}
\tightlist
\item
\end{enumerate}

\[
\operatorname{Tr} _ {\mathbf {M} \otimes \mathbf {M}} (a \otimes a ^ {\prime}) = \operatorname{Tr} _ {\mathbf {M}} (a) \operatorname{Tr} _ {\mathbf {M} ^ {\prime}} (a ^ {\prime})\tag{17}
\]

\[
\mathrm{N} _ {\mathbf {M} \otimes \mathbf {M} ^ {\prime}} (a \otimes a ^ {\prime}) = (\mathrm{N} _ {\mathbf {M}} (a)) ^ {n ^ {\prime}} (\mathrm{N} _ {\mathbf {M} ^ {\prime}} (a ^ {\prime})) ^ {n}.
\]

公式(16)由II，§4，第4号，公式(26)得出，而公式(17)由§8，第6号，公式(33)得出。

\subsection{代数中的范数与迹}

定义2. 设 A 是一个 K-代数，且它是有限维自由 K-模。对每个元素 \(a \in \mathbf{A}\) ，a 相对于 A 和 K 的迹（相应地，范数†，相应地，特征多项式）是 K-模 A 的自同态 \(x \mapsto ax\) 的迹（相应地，行列式，相应地，特征多项式）。

 \(a \in A\) 相对于 A 和 K 的迹、范数和特征多项式分别记为 \(\operatorname{Tr}_{\mathbf{A}/\mathbf{K}}(a)\) ， \(\operatorname{N}_{\mathbf{A}/\mathbf{K}}(a)\) 和 \(\operatorname{Pc}_{\mathbf{A}/\mathbf{K}}(a; \mathbf{X})\) ；当没有混淆风险时，我们在此记号中省略 K 甚至 A。注意， \(a \in A\) 的迹（相应地，范数、特征多项式）就是 a 相对于 A-模 \(A_{s}\) 的迹（相应地，范数，特征多项式）。

假设 A 是乘积 \(A_{1} \times A_{2} \times \cdots \times A_{m}\) ，即有限个 K 上的有限维代数（它们均为自由 K-模）的乘积。利用上述注记以及第 2 号的命题 1，对每个元素

\[
a = (a _ {1}, \dots , a _ {m}) \in \mathrm{A},
\]

\[
\mathrm{Tr} _ {\mathbf {A} / \mathbb {K}} (a) = \sum_ {i = 1} ^ {m} \mathrm{Tr} _ {\mathbf {A} _ {i} / \mathbb {K}} (a _ {i}), \qquad \mathrm{N} _ {\mathbf {A} / \mathbb {K}} (a) = \prod_ {i = 1} ^ {m} \mathrm{N} _ {\mathbf {A} _ {i} / \mathbb {K}} (a _ {i})\tag{18}
\]

\[
\mathrm{Pc} _ {\mathbf {A} / \mathbf {K}} (a; \mathbf {X}) = \prod_ {i = 1} ^ {m} \mathrm{Pc} _ {\mathbf {A} _ {i} / \mathbf {K}} (a _ {i}; \mathbf {X}).
\]

类似地，第 2 号的命题 2 表明，若 A 和 A' 是两个代数，它们都是自由 K-模，在 K 上的维数分别为 n 与 \(n'\) ，则对于 \(a \in A\) ， \(a' \in A'\) 。

\begin{enumerate}
\def\labelenumi{(\arabic{enumi})}
\setcounter{enumi}{18}
\tightlist
\item
\end{enumerate}

\[
\operatorname{Tr} _ {\mathbf {A} \otimes \mathbf {A} ^ {\prime}} (a \otimes a ^ {\prime}) = \operatorname{Tr} _ {\mathbf {A}} (a) \operatorname{Tr} _ {\mathbf {A} ^ {\prime}} (a ^ {\prime})\tag{20}
\]

\[
\mathrm{N} _ {\mathbf {A} \otimes \mathbf {A} ^ {\prime}} (a \otimes a ^ {\prime}) = (\mathrm{N} _ {\mathbf {A}} (a)) ^ {n ^ {\prime}} (\mathrm{N} _ {\mathbf {A} ^ {\prime}} (a ^ {\prime})) ^ {n}.
\]

最后，设 A 为 K 上的有限维代数，且它是自由 K-模，h 为 K 到交换环 \(K'\) 的一个同态， \(A' = A_{(K')}\) 为通过 h 作标量扩张而由 A 得到的 \(K'\) -代数。由第 1 号公式 (12) 可知，对所有 \(a \in A\) ，

\[
\begin{array}{r l} \mathrm{Tr} _ {\mathbf {A} ^ {\prime} / \mathbf {K} ^ {\prime}} (1 \otimes a) & = h (\mathrm{Tr} _ {\mathbf {A} / \mathbf {K}} (a)), \quad \mathrm{N} _ {\mathbf {A} ^ {\prime} / \mathbf {K} ^ {\prime}} (1 \otimes a) = h (\mathrm{N} _ {\mathbf {A} / \mathbf {K}} (a)) \\ \mathrm{Pc} _ {\mathbf {A} ^ {\prime} / \mathbf {K} ^ {\prime}} (1 \otimes a; \mathbf {X}) & = \tilde {h} (\mathrm{Pc} _ {\mathbf {A} / \mathbf {K}} (a; \mathbf {X})) \end{array}\tag{21}
\]

其中 \(\tilde{h}\) 是同态 \(\mathbf{K}[\mathbf{X}] \to \mathbf{K}'[\mathbf{X}]\) ，它由 \(h\) 导出。特别地，对于 \(\mathbf{K}' = \mathbf{K}[\mathbf{X}]\) ，记 \(\mathrm{A}[\mathbf{X}] = \mathrm{A} \otimes_{\mathbf{K}} \mathbf{K}[\mathbf{X}]\) ，我们得到

\[
\mathrm{Pc} _ {\mathbf {A} / \mathbf {K}} (a; \mathbf {X}) = \mathrm{N} _ {\mathbf {A} [ \mathbf {X} ] / \mathbf {K} [ \mathbf {X} ]} (\mathbf {X} - a).\tag{22}
\]

更一般地，若 \(\mathbf{K}'\) 是交换 \(\mathbf{K}\) -代数且 \(\mathbf{A}' = \mathbf{A} \otimes_{\mathbf{K}} \mathbf{K}'\) ，则对所有 \(x \in \mathbf{A}'\) ，

\[
\mathrm{Pc} _ {\mathbf {A} / \mathbf {K}} (a; x) = \mathrm{N} _ {\mathbf {A} ^ {\prime} / \mathbf {K} ^ {\prime}} (x - a).
\]

例子。(1) 设 A 是 K 上类型为 \((\alpha, \beta)\) 的二次代数， \((e_1, e_2)\) 是一个类型为 \((\alpha, \beta)\) 的基（§ 2，第 3 号）。对于 \(x = \xi e_1 + \eta e_2\) ， \(\operatorname{Tr}_{A/K}(x) = 2\xi + \beta \eta\) 和 \(\mathrm{N}_{A/K}(x) = \xi^2 + \beta \xi \eta - \alpha \eta^2\) ；因此这些函数与 \(x\) 的凯莱迹和凯莱范数相同（§ 2，第 24 号）。

\begin{enumerate}
\def\labelenumi{(\arabic{enumi})}
\setcounter{enumi}{1}
\item
  设 A 是 K 上的一个四元数代数。直接计算可以验证 \(\mathrm{Tr}_{\mathbf{A} / \mathbb{K}}(x) = 2\mathrm{T}(x)\) 和 \(\mathrm{N}_{\mathbf{A} / \mathbb{K}}(x) = (\mathrm{N}(x))^2\) ，其中 \(\mathbf{T}\) 和 \(\mathbf{N}\) 是凯莱迹和凯莱范数（§ 2，第 4 号）。
\item
  设 \(\mathbf{A} = \mathbf{M}_n(\mathbf{K})\) ，并令 \((E_{ij})\) （ \(\mathbf{A}\) 的典范基，II，§ 10，第 3 号）按字典序排列。则立即可见，对每个矩阵 \(\mathbf{X} = \sum_{i,j} \xi_{ij} E_{ij}\) ，（ \(n^2\) 阶的）自同态 \(Y \mapsto XY\) 的矩阵具有形式
\end{enumerate}

\[
\left( \begin{array}{c c c c} X & 0 & \ldots & 0 \\ 0 & X & \ldots & 0 \\ . & . & . & . \\ 0 & 0 & \ldots & X \end{array} \right)
\]

，因此 \(\operatorname{Tr}_{\mathbf{A} / \mathbb{K}}(X) = n.\operatorname{Tr}(X)\) 和 \(\mathrm{N}_{\mathbf{A} / \mathbb{K}}(X) = (\det (X))^n\)

\subsection{代数中范数与迹的性质}

命题3. 设 A 是一个具有有限基的 K-代数。元素 \(a \in \mathbf{A}\) 可逆的必要且充分条件是它的 \(\mathbf{N}_{\mathbf{A} / \mathbf{K}}(a)\) 在 \(\mathbf{K}\) 中可逆。

若 \(a\) 具有逆元 \(a'\) （位于 \(\mathbf{A}\) 中），则

\[
\mathrm{N} _ {\mathbf {A} / \mathbf {K}} (a) \mathrm{N} _ {\mathbf {A} / \mathbf {K}} (a ^ {\prime}) = \mathrm{N} _ {\mathbf {A} / \mathbf {K}} (a a ^ {\prime}) = \mathrm{N} _ {\mathbf {A} / \mathbf {K}} (1) = 1
\]

由第 1 号公式 (3)。反之，若 \(\mathbf{N}_{\mathbf{A} / \mathbb{K}}(a)\) 可逆，则自同态 \(h: x \mapsto ax\) 是双射（§ 8，第 2 号，定理 1）。于是存在 \(a' \in A\) 使得 \(aa' = 1\) ；于是 \(h(a'a - 1) = aa'a - a = (aa' - 1)a = 0\) ，从而 \(aa' = 1\) ，因为 \(h\) 是单射。因此 \(a'\) 是 \(a\) 的逆元。

命题4. 设 A 是一个具有有限基的 K-代数。对所有 \(a \in \mathbf{A}\) ， \(\operatorname{Pc}_{\mathbf{A} / \mathbf{K}}(a; a) = 0\) 。

这直接由凯莱-哈密顿定理（§ 8，第 11 号，命题 20）得出。

命题5. 设 A 是一个 K-代数，m 是 A 的一个双边理想。假设 \(A_{0} = A/m\) 是维数为 n 的自由 K-模，存在整数 r \textgreater{} 0 使得 \(m^{r} = \{0\}\) ，并且 \(m^{i-1}/m^{i}\) 是自由 \(A_{0}\) -模，其维数为 \(s_{i}\) （ \(1 \leqslant i \leqslant r\) ）。设 \(s = s_{1} + \cdots + s_{r}\) ，且对所有 \(a \in A\) ，设 \(a_{0}\) 表示 a 模 m 所在的类。则 A 是维数为 ns 的自由 K-模，且对所有 \(a \in A\) ，

\[
\operatorname{Tr} _ {\mathbf {A}} (a) = s. \operatorname{Tr} _ {\mathbf {A} _ {0}} (a _ {0}), \quad \mathrm{N} _ {\mathbf {A}} (a) = (\mathrm{N} _ {\mathbf {A} _ {0}} (a _ {0})) ^ {s}\tag{23}
\]

\[
\mathrm{Pc} _ {\mathbf {A}} (a; \mathbf {X}) = (\mathrm{Pc} _ {\mathbf {A} _ {0}} (a _ {0}; \mathbf {X})) ^ {s}.
\]

根据 II，§1，第 13 号，命题 25， \(\mathfrak{m}^{i - 1} / \mathfrak{m}^{i}\) 是维数为 \(ns_i\) 的自由 K-模。因此，第 2 号中的命题 1 可以应用于 \(\mathbf{P}_i = \mathfrak{m}^{i - 1} / \mathfrak{m}^{i}\) ；

这首先表明 A 是一个维数为 \(n(s_{1} + \cdots + s_{r}) = ns\) 的自由 K-模。此外，该假设蕴含 A-模 \(P_{i}\) 同构于 \(s_{i}\) 个子模的直和，这些子模均同构于 A-模 \(A_{0}\) ；由第 2 号命题 1，因此 \(\mathrm{N}_{\mathrm{P}_{i}}(a) = \mathrm{N}_{\mathrm{A}_{0}}(a)^{s_{i}}\) ；最终因此

\[
\mathrm{N} _ {\mathbf {A}} (a) = \mathrm{N} _ {\mathbf {A} _ {0}} (a) ^ {s}.
\]

在此公式中， \(\mathrm{N}_{\mathbf{A}_{0}}(a)\) 是通过将 \(A_{0}\) 视为左 A-模来定义的，它等于 K-线性映射 \(x \mapsto ax\) 在 \(A_{0}\) 到自身上的行列式；但由于 \(ax = a_{0}x\) 对 \(x \in A_{0}\) 成立，故 \(\mathrm{N}_{\mathbf{A}_{0}}(a) = \mathrm{N}_{\mathbf{A}_{0}}(a_{0})\) ，这就完成了范数公式 (23) 的证明。另外两个公式可类似地证明。

命题6. 设 A 是 K 上的一个交换代数，且在 K 上具有有限基，V 是一个 A-模，且在 A 上具有有限基。则 V 在 K 上具有有限基，并且对 V 的每个 A-自同态 u，若 \(u_{\mathbf{K}}\) 是将u视为V的K自同态的映射，

\[
\operatorname{Tr} (u _ {\mathbb {K}}) = \operatorname{Tr} _ {\mathbf {A} / \mathbb {K}} (\operatorname{Tr} (u)), \quad \det (u _ {\mathbb {K}}) = \mathrm{N} _ {\mathbf {A} / \mathbb {K}} (\det (u))\tag{24}
\]

\[
\operatorname{Pc} (u _ {\mathbf {K}}; \mathbf {X}) = \mathrm{N} _ {\mathbf {A} [ \mathbf {X} ] / \mathbf {K} [ \mathbf {X} ]} (\operatorname{Pc} (u; \mathbf {X})).
\]

设 \((a_{i})_{1\leqslant i\leqslant m}\) 是 A 在 K 上的一个基， \((e_{j})_{1\leqslant j\leqslant n}\) 是 V 在 A 上的一个基；则 \((a_{i}e_{j})\) 是 V 在 K 上的一个基（II，§1，第13号，命题25）。另一方面，将第二个公式应用于自同态 \(\mathbf{X} - \bar{\boldsymbol{u}}\) （即 A{[}X{]}-模 A{[}X{]} \(\otimes_{\mathbf{A}}\mathrm{V}\) 的自同态）（§ 8，第 10 号），即可导出公式 (24) 中的第三个。因此，只需证明 (24) 中的前两个公式即可。我们首先建立以下引理：

引理1. 设 \(\mathbf{X}_{ij}\) ( \(1 \leqslant i \leqslant n\) , \(1 \leqslant j \leqslant n\) ) 是 \(n^2\) 个未定元， \(X\) 是元素为 \((\mathbf{X}_{ij})\) 的 \(n\) 阶方阵， \(\mathbf{D}(\mathbf{X}_{11}, \ldots, \mathbf{X}_{nn}) \in \mathbf{Z}[\mathbf{X}_{11}, \ldots, \mathbf{X}_{nn}]\) 是行列式 \(\det(X)\) 。另一方面，设 A 为交换环， \(\mathbf{M}_{ij}\) ( \(1 \leqslant i \leqslant n\) , \(1 \leqslant j \leqslant n\) ) 是 A 上的 \(n^2\) 个 \(m\) 阶矩阵，它们两两可交换，M 是 A 上的 \(mn\) 阶方阵，它可以表示为矩阵的分块方阵（II，§10，第7号）

\[
M = \left( \begin{array}{c c c c} M _ {1 1} & M _ {1 2} & \ldots & M _ {1 n} \\ M _ {2 1} & M _ {2 2} & \ldots & M _ {2 n} \\ \cdot & \cdot & \cdot & \cdot \\ M _ {n 1} & M _ {n 2} & \ldots & M _ {n n} \end{array} \right)
\]

那么M的行列式等于方阵的行列式

\[
\mathrm{D} (M _ {1 1}, \dots , M _ {n n})
\]

该方阵的阶为 m。

我们对 \(n\) 进行归纳，情形 \(n = 0\) 和 \(n = 1\) 是平凡的。

设 \(Z\) 为一个新的未定元， \(N_{ij}\) 表示矩阵 \(M_{ij} + \delta_{ij}ZI_m\) ( \(\delta_{ij}\) 为克罗内克指标)。若 \(D^{ij}(X_{11},\ldots,X_{nn})\) 是 \(X_{ij}\) 在矩阵 \(X\) （§ 8，第 6 号）中的代数余子式，则

\[
\mathrm{X} _ {j l} \mathrm{D} ^ {k l} (\mathrm{X} _ {1 1}, \dots , \mathrm{X} _ {n n}) = \delta_ {j k} \mathrm{D} (\mathrm{X} _ {1 1}, \dots , \mathrm{X} _ {n n})\tag{25}
\]

（§ 8，第 6 号，公式 (28)）。我们记 \(N_{ij}' = \mathbf{D}^{ij}(N_{11}, \ldots, N_{nn})\) ，这是一个 \(m\) 阶方阵，位于 \(\mathrm{A}[Z]\) 上；并考虑乘积 \(N.U\) ，其中

\[
U = \left( \begin{array}{c c c c} N _ {1 1} ^ {\prime} & 0 & \ldots & 0 \\ N _ {1 2} ^ {\prime} & I _ {m} & \ldots & 0 \\ \cdot & \cdot & \cdot & \cdot \\ N _ {1 n} ^ {\prime} & 0 & \ldots & I _ {m} \end{array} \right),
\]

\[
N = \left( \begin{array}{c c c c} N _ {1 1} & N _ {1 2} & \ldots & N _ {1 n} \\ N _ {2 1} & N _ {2 2} & \ldots & N _ {2 n} \\ \cdot & \cdot & \cdot & \cdot \\ N _ {n 1} & N _ {n 2} & \ldots & N _ {n n} \end{array} \right)
\]

按分块进行此乘积（II，§ 10，第 5 号）并利用公式 (25)，我们得到

\[
N. U = \left( \begin{array}{c c c c} P & N _ {1 2} & \ldots & N _ {1 n} \\ 0 & N _ {2 2} & \ldots & N _ {2 n} \\ \cdot & \cdot & \cdot & \cdot \\ 0 & N _ {n 2} & \ldots & N _ {n n} \end{array} \right)
\]

其中我们记作 \(P = \mathrm{D}(N_{11},\ldots ,N_{nn})\) 。设

\[
Q = \left( \begin{array}{c c c c} N _ {2 2} & \ldots & N _ {2 n} \\ \cdot & \cdot & \cdot & \cdot \\ N _ {n 2} & \ldots & N _ {n n} \end{array} \right)
\]

，这是一个 \(m(n - 1)\) 阶矩阵；则（§8，第6号，公式 (31)）(det \(N\) ) ( \(\det U\) ) = ( \(\det P\) ) ( \(\det Q\) ) 且 \(\det U = \det N_{11}'\) 。但根据归纳假设， \(\det Q = \det(\mathrm{D}^{11}(N_{11},\ldots ,N_{nn})) = \det N_{11}'\) ，并且根据 \(N_{ij}\) 的定义，显然 \(\det Q\) 是 A{[}Z{]} 中一个次数为 \(m(n - 1)\) 的多项式，其 \(Z^{m(n - 1)}\) 项的系数为 1；由此立即得出 \(\det Q\) 在分次代数 A{[}Z{]} 中不是零因子。因此我们得出结论： \(\det N = \det(\mathrm{D}(N_{11},\ldots ,N_{nn}))\) 在 A{[}Z{]} 中成立；如果我们在这些多项式中将 Z 替换为 0，那么 \(\det M = \det(\mathrm{D}(M_{11},\ldots ,M_{nn}))\) 。

证明了这一引理之后，K-模 V 是诸 K-模 \(Ae_{j}\) ( \(1 \leqslant j \leqslant n\) ) 的直和；我们记 \(u(e_{j}) = \sum_{k=1}^{n} c_{jk} e_{k}\) 。对每个元素 \(xe_{j} \in Ae_{j}\) ，其中 \(x \in A\) ， \(u(xe_{j})\) 在 \(Ae_{k}\) 中的分量是 \(xc_{jk} e_{k}\) ；由此可知， \(u_{k}\) 关于 K-模 V 的基 \((a_{i} e_{j})\) 的矩阵可以表示为分块方阵 \((M_{jk})\) 的形式，其中 \(M_{jk}\) 是 K-线性映射 \(xe_{j} \mapsto xc_{jk} e_{k}\) （它把 \(Ae_{j}\) 映到 \(Ae_{k}\) ）关于这两个 K-模的基 \((a_{i} e_{j})_{1 \leqslant i \leqslant m}\) 和 \((a_{i} e_{k})_{1 \leqslant i \leqslant m}\) （II，§ 10，第 5 号）的矩阵。若对所有 \(t \in A\) ， \(M(t)\) 表示 A 在 K 上的基 \((a_{i})_{1 \leqslant i \leqslant m}\) 下 A 的自同态 \(x \mapsto xt\) 的矩阵，则 \(M_{jk} = M(c_{jk})\) ；由于 \(t \mapsto M(t)\) 是一个环同态，矩阵 \(M_{jk}\) 两两可交换。则

\[
\det u _ {\mathrm{K}} = \det (\mathrm{D} (M _ {1 1}, \dots , M _ {n n}))
\]

由引理1。但由于 \(t \mapsto M(t)\) 是一个环同态， \(\mathrm{D}(M_{11},\ldots ,M_{nn})\) 是 A 的 K-自同态 \(x \mapsto x.\det (c_{jk})\) 关于基 \((a_t)\) 的矩阵；根据定义，其行列式因此为 \(\mathrm{N}_{\mathbf{A} / \mathbf{K}}(\det (u))\) ，这证明了 (24) 中的第二个公式。另一方面，

\[
\operatorname{Tr} \left(u _ {\mathrm{K}}\right) = \sum_ {j = 1} ^ {n} \operatorname{Tr} \left(M _ {j j}\right) = \sum_ {j = 1} ^ {n} \operatorname{Tr} _ {\mathrm{A} / \mathrm{K}} \left(c _ {j j}\right) = \operatorname{Tr} _ {\mathrm{A} / \mathrm{K}} \left(\sum_ {j = 1} ^ {n} c _ {j j}\right) = \operatorname{Tr} _ {\mathrm{A} / \mathrm{K}} (\operatorname{Tr} (u))
\]

并且命题6的证明完成。

推论。设A是交换K-代数，在K上具有有限基，B是A-代数，在A上具有有限基。则B在K上具有有限基，并且，对所有 \(b \in B\) （“传递性公式”）

\[
\operatorname{Tr} _ {\mathbf {B} / \mathbf {K}} (b) = \operatorname{Tr} _ {\mathbf {A} / \mathbf {K}} (\operatorname{Tr} _ {\mathbf {B} / \mathbf {A}} (b)), \quad \mathrm{N} _ {\mathbf {B} / \mathbf {K}} (b) = \mathrm{N} _ {\mathbf {A} / \mathbf {K}} (\mathrm{N} _ {\mathbf {B} / \mathbf {A}} (b))\tag{26}
\]

\[
\mathrm{Pc} _ {\mathbf {B} / \mathbf {K}} (b; \mathbf {X}) = \mathrm{N} _ {\mathbf {A} [ \mathbf {X} ] / \mathbf {K} [ \mathbf {X} ]} (\mathrm{Pc} _ {\mathbf {B} / \mathbf {A}} (b; \mathbf {X})).
\]

这直接从命题 6 得出，取 \(\mathbf{V} = \mathbf{B}\) 和 \(u(x) = bx\) 。

注记。假设 K 到 A 的同态 \(\lambda \mapsto \lambda, 1\) 是单射，并将 K 与其在 A 中的像等同；假设 A 作为 K-模具有有限基 \((e_i)_{1 \leqslant i \leqslant n}\) 。设 \(s\) 是 A 的一个自同构，使得 \(s(\mathbf{K}) = \mathbf{K}\) 。设 \(a\) 是 A 的一个元素；则通过转移结构

\begin{enumerate}
\def\labelenumi{(\arabic{enumi})}
\setcounter{enumi}{26}
\tightlist
\item
\end{enumerate}

\[
\operatorname{Tr} _ {\mathbf {A} / \mathbf {K}} (s (a)) = s (\operatorname{Tr} _ {\mathbf {A} / \mathbf {K}} (a))\tag{28}
\]

\[
\mathrm{N} _ {\mathbf {A} / \mathbf {K}} (s (a)) = s (\mathrm{N} _ {\mathbf {A} / \mathbf {K}} (a)).
\]

同时考虑 A 的一个导子 D（§ 10，第 2 号），使得 D(K) ⊂ K，并记 \(D (e _ {i}) = \sum_ {j = 1} ^ {n} e _ {j} \mu_ {j i}\) ，其中 \(\mu_ {j i} \in \mathbf {K}\) ；记

\[
a e _ {i} = \sum_ {j = 1} ^ {n} e _ {j} \lambda_ {j i} \quad \text { with } \quad \lambda_ {j i} \in \mathbf {K}.
\]

那么

\[
\mathrm{D} (a) e _ {i} + a \mathrm{D} (e _ {i}) = \mathrm{D} (a e _ {i}) = \sum_ {j = 1} ^ {n} (\mathrm{D} (e _ {j}) \lambda_ {j i} + e _ {j} \mathrm{D} (\lambda_ {j i})).
\]

由此可得

\[
\mathbf {D} (a) e _ {i} = \sum_ {j = 1} ^ {n} e _ {j} \nu_ {j i}
\]

其中 \(\nu_{ji} = \mathrm{D}(\lambda_{ji}) + \sum_{k=1}^{n} (\mu_{jk}\lambda_{ki} - \lambda_{jk}\mu_{ki})\) 。由于 \(\sum_{i,k} (\mu_{ik}\lambda_{ki} - \lambda_{ik}\mu_{ki}) = 0\) ，因此

 \(\mathrm{Tr}_{\mathbf{A} / \mathbb{K}}(\mathbf{D}(a)) = \sum_{i = 1}^{n}\mathbf{D}(\lambda_{ii})\) ，换言之

\[
\operatorname{Tr} _ {\mathbf {A} / \mathbf {K}} (\mathbf {D} (a)) = \mathbf {D} (\operatorname{Tr} _ {\mathbf {A} / \mathbf {K}} (a)). *\tag{29}
\]

\subsection{代数的判别式}

定义 3. 设 A 是一个 K-代数，具有由 n 个元素构成的有限基。A 中 n 个元素的序列 \((x_{1},\ldots,x_{n})\) 关于 K 的判别式，记作 \(\mathrm{D}_{\mathrm{A}/\mathrm{K}}(x_{1},\ldots,x_{n})\) ，定义为下述方阵的判别式

\[
\left(\operatorname{Tr} _ {\mathbf {A} / \mathbb {K}} \left(x _ {i} x _ {j}\right)\right) _ {1 \leqslant i \leqslant n, 1 \leqslant j \leqslant n}.
\]

首先考虑 A 在 K 上的一个基 \((e_i)_{1\leqslant i\leqslant n}\) ，并记

\[
e _ {i} e _ {j} = \sum_ {k = 1} ^ {n} c _ {i j k} e _ {k} \quad \text { with } \quad c _ {i j k} \in \mathbf {K}.\tag{30}
\]

那么 \(\mathrm{Tr}_{\mathbf{A} / \mathbf{K}}(e_i) = \sum_{s=1}^{n} c_{iss}\) ，由此 \(\mathrm{Tr}_{\mathbf{A} / \mathbf{K}}(e_i e_j) = \sum_{k,s} c_{ijk} c_{kss}\) ，进而

\[
\mathrm{D} _ {\mathrm{A} / \mathrm{K}} (e _ {1}, \dots , e _ {n}) = \det \left(\left(\sum_ {k, s} c _ {i j k} c _ {k s s}\right) _ {1 \leqslant i \leqslant n, 1 \leqslant j \leqslant n}\right).\tag{31}
\]

现在令 \((x_{i})_{1\leqslant i\leqslant n},(x_{i}^{\prime})_{1\leqslant i\leqslant n}\) 是 A 的 \(n\) 个元素的两个序列，并假设存在一个 \(n\) 阶方阵 \(M = (m_{ij})\) ，其系数在 K 中，使得 \(x_{i} = \sum_{j=1}^{n} m_{ij} x_{j}'\) 对 \(1 \leqslant i \leqslant n\) 成立。我们记

\[
T = \left(\operatorname{Tr} _ {\mathrm{A} / \mathrm{K}} \left(x _ {i} x _ {j}\right)\right) _ {1 \leqslant i \leqslant n, 1 \leqslant j \leqslant n}, \quad T ^ {\prime} = \left(\operatorname{Tr} _ {\mathrm{A} / \mathrm{K}} \left(x _ {i} ^ {\prime} x _ {j} ^ {\prime}\right)\right) _ {1 \leqslant i \leqslant n, 1 \leqslant j \leqslant n}.
\]

那么 \(\mathrm{Tr}_{\mathbf{A} / \mathbf{K}}(x_i x_j) = \sum_{p,q} m_{ip} m_{jq} \mathrm{Tr}_{\mathbf{A} / \mathbf{K}}(x_p' x_q')\) ，因此 \(T = M. T'.^t M\) ；于是行列式的乘法规则给出

\[
\det T = \det M. \det T ^ {\prime}. \det ^ {t} M = (\det M) ^ {2} \det T ^ {\prime}
\]

最终得到

\[
\mathrm{D} _ {\mathrm{A} / \mathrm{K}} (x _ {1}, \dots , x _ {n}) = (\det M) ^ {2} \mathrm{D} _ {\mathrm{A} / \mathrm{K}} (x _ {i} ^ {\prime}, \dots , x _ {n} ^ {\prime}).\tag{32}
\]

上述公式特别表明，A 在 K 上的两组基的判别式相差 K 中一个可逆元素的平方，因此生成 K 的同一个（主）理想。这个理想 \(\Delta_{A/K}\) 称为 A 在 K 上的判别式理想；由公式 (32)，A 中仅项的顺序不同的任意 n 个元素的序列都具有相同的判别式，因为置换矩阵的行列式等于 \(\pm1\) 。

例子。(1) 若 A 是 K 上类型为 \((\alpha, \beta)\) 的二次代数，则（采用 § 2 第 3 号的记号） \(\operatorname{Tr}(e_1) = 2\) ， \(\operatorname{Tr}(e_2) = \beta\) ，

\[
\operatorname{Tr} (e _ {2} ^ {2}) = \alpha \operatorname{Tr} (e _ {1}) + \beta \operatorname{Tr} (e _ {2}) = 2 \alpha + \beta^ {2},
\]

由此 \(\mathbf{D}_{\mathbf{A} / \mathbb{K}}(e_1,e_2) = \beta^2 +4\alpha .\)

\begin{enumerate}
\def\labelenumi{(\arabic{enumi})}
\setcounter{enumi}{1}
\item
  设 \(\mathbf{A} = \mathbf{K}[\mathbf{X}] / \mathbf{K}[\mathbf{X}]\mathbf{P}\) ，其中 \(\mathrm{P}(\mathbf{X}) = \mathbf{X}^3 + p\mathbf{X} + q\) ，于是若 \(x\) 是 \(\mathbf{X}\) 在 \(\mathbf{A}, 1, x, x^2\) 中的像，则它们构成 \(\mathbf{A}\) 在 \(\mathbf{K}\) 上的一组基，且 \(x^3 = -px - q\) 。立即可以看出 \(\operatorname{Tr}(1) = 3\) ， \(\operatorname{Tr}(x) = 0\) ， \(\operatorname{Tr}(x^2) = -2p\) ；利用关系 \(x^3 = -px - q\) ， \(\operatorname{Tr}(x^3) = -3q\) ， \(\operatorname{Tr}(x^4) = 2p^2\) ，由此容易得到 \(\mathrm{D}_{\mathbf{A} / \mathbf{K}}(1, x, x^2) = -4p^3 - 27q^2\) 。
\item
  设 A 是 K 上类型为 \((\alpha, \beta, \gamma)\) 的四元数代数， \((1, i, j, k)\) 是 A 的一个类型为 \((\alpha, \beta, \gamma)\) 的基；利用 § 3 第 5 号公式 (30)，容易得出 \(\operatorname{Tr}(1) = 4\) ， \(\operatorname{Tr}(i) = 2\beta\) ， \(\operatorname{Tr}(j) = \operatorname{Tr}(k) = 0\) ，于是
\end{enumerate}

\[
\mathrm{D} _ {\mathrm{A/K}} (1, i, j, k) = - 16 \gamma^ {2} (\beta^ {2} + 4 \alpha) ^ {2}.
\]

\begin{enumerate}
\def\labelenumi{(\arabic{enumi})}
\setcounter{enumi}{3}
\tightlist
\item
  设 \(\mathbf{A} = \mathbf{M}_n(\mathbf{K})\) ，并考虑典范基 \((E_{ij})_{1\leqslant i\leqslant n,1\leqslant j\leqslant n}\) （ \(\mathbf{A}\) 在 \(\mathbf{K}\) 上的）(II, § 10, 第 3 号)。立即可见 \(\operatorname{Tr}_{\mathbf{A} / \mathbf{K}}(E_{ij}) = 0\) （若 \(j\neq i\) ），且 \(\operatorname{Tr}_{\mathbf{A} / \mathbf{K}}(E_{ii}) = n\) 对所有 \(i\) 成立；由此不难得出，矩阵 \((\operatorname{Tr}(E_{ij}E_{hk}))\) （ \(n^2\) 阶）具有 \(n.P\) 的形式，其中 \(P\) 是一个置换矩阵，从而 \(\mathrm{D}_{\mathbf{A} / \mathbf{K}}((E_{ij})) = \pm n^{n^2}\) 。
\end{enumerate}

\section{导子}

在本节中，除非另有说明，所考虑的代数不假定为结合的，也不一定具有单位元；K 表示一个交换环。

\subsection{交换因子}

在本段中，当我们不加说明地提及分次时，总是指类型为 \(\Delta\) 的分次，其中 \(\Delta\) 是一个以加法记写的交换群。在本段中， \(\Delta\) 上取值于 Z 的交换因子称为 \(\Delta\) 上的交换因子（ \(\S\) 4，第 7 号，定义 6）。因此， \(\Delta\) 上的一个交换因子等同于一个映射 \(\varepsilon\colon(\alpha,\beta)\mapsto\varepsilon_{\alpha\beta}=\varepsilon(\alpha,\beta)\) ，它把 \(\Delta\times\Delta\) 映到乘法群 \(\{-1,1\}\) 中，使得对 \(\alpha,\alpha^{\prime},\beta,\beta^{\prime}\) （均为 \(\Delta\) 中的元素），

\[
\left\{ \begin{array}{l} \varepsilon (\alpha + \alpha^ {\prime}, \beta) = \varepsilon (\alpha , \beta) \varepsilon (\alpha^ {\prime}, \beta) \\ \varepsilon (\alpha , \beta + \beta^ {\prime}) = \varepsilon (\alpha , \beta) \varepsilon (\alpha , \beta^ {\prime}) \\ \varepsilon (\beta , \alpha) = \varepsilon (\alpha , \beta). \end{array} \right.\tag{1}
\]

由此可得 \(\varepsilon(2\alpha, \beta) = \varepsilon(\alpha, 2\beta) = 1\) .

当 \(\Delta = \mathbf{Z}\) 时，每个交换因子 \(\varepsilon\) 由 \(\varepsilon(1, 1)\) 的值决定；因此这样的因子只有两个，第一个由

\[
\varepsilon (p, q) = 1 \quad \text { for } \quad p, q \text { in } \mathbf {Z}\tag{2}
\]

定义，第二个由

\[
\varepsilon (p, q) = (- 1) ^ {p q} \quad \text { for } \quad p, q \text { in } \mathbf {Z}.\tag{3}
\]

\subsection{导子的一般定义}

考虑一个交换环 K，六个类型为 \(\Delta\) 的分次 K-模：A, A', A'', B, B', B''，以及三个 K-线性映射

\[
\mu \colon \mathrm{A} \times \mathrm{A} ^ {\prime} \rightarrow \mathrm{A} ^ {\prime \prime}, \quad \lambda_ {1} \colon \mathrm{B} \times \mathrm{A} ^ {\prime} \rightarrow \mathrm{B} ^ {\prime \prime}, \quad \lambda_ {2} \colon \mathrm{A} \times \mathrm{B} ^ {\prime} \rightarrow \mathrm{B} ^ {\prime \prime}
\]

使得相应的 K-线性映射

\[
\mathrm{A} \otimes_ {\mathrm{K}} \mathrm{A} ^ {\prime} \rightarrow \mathrm{A} ^ {\prime \prime}, \quad \mathrm{B} \otimes_ {\mathrm{K}} \mathrm{A} ^ {\prime} \rightarrow \mathrm{B} ^ {\prime \prime}, \quad \mathrm{A} \otimes_ {\mathrm{K}} \mathrm{B} ^ {\prime} \rightarrow \mathrm{B} ^ {\prime \prime}
\]

是次数为 0 的分次映射。像 \(\mu(a, a')\) （对于 \(a \in A\) ， \(a' \in A'\) ）简记为 \(a, a'\) 甚至 \(aa'\) ，另外两个双线性映射同理。因此 \(a, a'\) 的次数是 \(a\) 与 \(a'\) 的次数之和。

定义 1. 在上述设定和一个交换因子 \(\varepsilon\) （ \(\Delta \times \Delta\) 上的）之下，一个 \(\varepsilon\) -导子（或 \((\mathbf{K}, \varepsilon)\) -导子）（次数为 \(\delta \in \Delta\) ，从 \((\mathbf{A}, \mathbf{A}', \mathbf{A}'')\) 到 \((\mathbf{B}, \mathbf{B}', \mathbf{B}'')\) ）指的是次数为 \(\delta\) 的分次 K-线性映射的一个三元组：

\[
d \colon \mathrm{A} \rightarrow \mathrm{B}, \quad d ^ {\prime} \colon \mathrm{A} ^ {\prime} \rightarrow \mathrm{B} ^ {\prime}, \quad d ^ {\prime \prime} \colon \mathrm{A} ^ {\prime \prime} \rightarrow \mathrm{B} ^ {\prime \prime}
\]

的分次K-线性映射三元组，使得对每个齐次元素 \(a \in A\) 以及每个元素 \(a' \in A'\)

\[
d ^ {\prime \prime} (a. a ^ {\prime}) = (d a). a ^ {\prime} + \varepsilon_ {\delta , \deg (a)} a. (d ^ {\prime} a ^ {\prime}).\tag{4}
\]

显然，由线性性，只需在 \(a\) 和 \(a'\) 分别遍历 A 和 A' 的生成系时验证关系 (4) 即可。

把这三个映射 \(d, d', d''\) 用同一个字母 \(d\) 来表示往往是方便的（其合理性在于同样用 \(d\) 表示次数为 \(\delta\) 的分次 K-线性映射

\[
(a, a ^ {\prime}, a ^ {\prime \prime}) \mapsto (d a, d ^ {\prime} a ^ {\prime}, d ^ {\prime \prime} a ^ {\prime \prime})
\]

它把 \(\mathbf{A} \oplus \mathbf{A}' \oplus \mathbf{A}''\) 映到 \(\mathbf{B} \oplus \mathbf{B}' \oplus \mathbf{B}''\) 中。于是关系 (4) 可以更简单地写成

\[
d (a. a ^ {\prime}) = (d a). a ^ {\prime} + \varepsilon_ {\delta . \deg (a)} a. (d a ^ {\prime}).\tag{5}
\]

从 (A, A', A'') 到 (B, B', B'') 的给定次数的 \(\varepsilon\) -导子构成分次线性映射 K-模的一个子 K-模

\[
\operatorname{Homgr} _ {\mathbb {K}} (\mathrm{A} \oplus \mathrm{A} ^ {\prime} \oplus \mathrm{A} ^ {\prime \prime}, \mathrm{B} \oplus \mathrm{B} ^ {\prime} \oplus \mathrm{B} ^ {\prime \prime}).
\]

当 \(\varepsilon(\alpha, \beta) = 1\) 对所有 \(\alpha, \beta\) （ \(\Delta\) 中的）成立时，我们简称导子（或 K-导子）而不称 \(\varepsilon\) -导子。导子构成下列 K-模的一个子 K-模

\[
\operatorname{Hom} _ {\mathbb {K}} (\mathrm{A} \oplus \mathrm{A} ^ {\prime} \oplus \mathrm{A} ^ {\prime \prime}, \mathrm{B} \oplus \mathrm{B} ^ {\prime} \oplus \mathrm{B} ^ {\prime \prime}).
\]

当 \(\Delta = Z\) 且 \(\varepsilon(p, q) = (-1)^{pq}\) 时，每个偶次数的 \(\varepsilon\) -导子都是导子；每个奇次数的 \(\varepsilon\) -导子通常称为反导子（或 K-反导子）；于是反导子 d 满足

\[
d (a. a ^ {\prime}) = (d a). a ^ {\prime} + (- 1) ^ {\deg (a)} a. (d a ^ {\prime})\tag{6}
\]

对齐次元素 \(a \in A\) 成立。

注记。(1) 导子的概念可以推广到非分次模上，只需约定赋予这些模平凡分次结构即可。

\begin{enumerate}
\def\labelenumi{(\arabic{enumi})}
\setcounter{enumi}{1}
\tightlist
\item
  如果只考虑 \(\varepsilon\) -导子中次数为给定值 \(\delta\) 的那些，交换因子 \(\varepsilon\) 可以按如下方式消去：把双线性映射 \(\lambda_2: \mathbf{A} \times \mathbf{B}' \to \mathbf{B}''\) 替换为双线性映射 \(\lambda_2': \mathbf{A} \times \mathbf{B}' \to \mathbf{B}''\) ，使得对每个齐次元素 \(a\) （ \(\mathbf{A}\) 中的）以及所有 \(b' \in \mathbf{B}'\) ，
\end{enumerate}

\[
\lambda_ {2} ^ {\prime} (a, b ^ {\prime}) = \varepsilon_ {\delta , \deg (a)} \lambda_ {2} (a, b ^ {\prime}).
\]

那么 \(d\) 是关于双线性映射 \(\mu, \lambda_{1}, \lambda_{2}^{\prime}\) 的导子。

上述给出的 \(\varepsilon\) -导子的一般定义特别用于两种情况：

情形（I）：A = B， \(A' = B'\) , \(A'' = B''\) 且三个双线性映射 \(\mu\) , \(\lambda_{1}\) , \(\lambda_{2}\) 等于同一个映射。

情形（II）： \(\mathbf{A} = \mathbf{A}' = \mathbf{A}''\) ， \(\mathbf{B} = \mathbf{B}' = \mathbf{B}''\) ，因此（对于 \(\mu\) ） \(\mathbf{A}\) 是一个分次代数，且这两个 K-双线性映射

\[
\lambda_ {1} \colon \mathrm{B} \times \mathrm{A} \rightarrow \mathrm{B}, \quad \lambda_ {2} \colon \mathrm{A} \times \mathrm{B} \rightarrow \mathrm{B}\tag{7}
\]

使得相应的 K-线性映射 \(B \otimes_{K} A \to B\) ， \(A \otimes_{K} B \to B\) 是次数为 0 的分次映射。此时，一个 \(\varepsilon\) -导子（次数为 \(\delta\) ，从 A 到 B）是指一个分次 K-线性映射 \(d: A \to B\) （次数为 \(\delta\) ），使得对 A 中每个齐次元素 x 以及每个 \(y \in A\) ，我们有关系式

\[
d (x y) = (d x) y + \varepsilon_ {\delta , \deg (a)} x (d y).\tag{8}
\]

特别地，在情形 (II) 中考虑这样的情形：A 是带单位元的结合 K-代数，且 \(\lambda_{1}\) 和 \(\lambda_{2}\) 是一个 (A, A)-双模（§ 4，第 3 号，定义 3）的外部法则。特别当 A 和 B 是两个带单位元的结合 K-代数、给定分次 K-代数的一个带单位元的同态 \(\rho: A \to B\) 并在 B 上考虑由这两个外部法则定义的 (A, A)-双模结构时，上述情形成立

\[
\lambda_ {1} \colon (b, a) \mapsto b \rho (a), \quad \lambda_ {2} \colon (a, b) \mapsto \rho (a) b
\]

对于 \(a\in \mathbf{A},b\in \mathbf{B}.\)

情形（I）和（II）具有以下共同情形：考虑一个分次 K-代数 A，取 B = A，映射 (7) 均为 A 上的乘法。此时我们谈论分次代数 A 的 ε-导子（或 (K, ε)-导子）：它是 A 到自身的、次数为 δ 的分次 K-线性映射，对 A 中每个齐次元素 x 以及所有 \(y \in A\) 满足 (8)。特别地，若 A 是一个分次环（视为（结合的）Z-代数），我们就谈论环 A 的 ε-导子。

设A是一个含单位元的交换结合K-代数，B是一个A-模；当我们谈到从A到B的导子时，总是指B上由A-模结构导出的A-双模结构；此时公式

\[
d (x y) = x (d y) + y (d x) \quad \text { for } \quad x \in \mathrm{A}, y \in \mathrm{A}\tag{9}
\]

对这样的导子 \(d: A \rightarrow B\) 成立。

\subsection{导子的例子}

例 1. *设 A 是 R 到 R 的可微映射构成的 R-代数，并设 \(x_0\) 是 R 的一个点；R 可视为一个 A-模，其外运算为 \((f, a) \mapsto f(x_0)a\) 。则映射 \(f \mapsto \mathrm{D}f(x_0)\) 是一个导子，因为（实变函数，I，§ 1，第 3 号）

\[
(\mathbf {D} (f g)) (x _ {0}) = (\mathbf {D} f (x _ {0})) g (x _ {0}) + f (x _ {0}) (\mathbf {D} g (x _ {0})). *
\]

例 2. *设 X 是一个 C∞ 类可微流形，并设 A 为 X 上微分形式的分次 R-代数。将每个微分形式 ω 映到其外微分 dω 的映射是次数为 +1 的反导子（可微与解析流形，R，§ 8）。*

例 3. 设 A 是一个结合 K-代数。对所有 \(a \in \mathbf{A}\) ，映射 \(x \mapsto ax - xa\) 是代数 A 的一个导子（参见第 6 号）。

例 4. 设 M 是一个 K-模，A 是带通常分次的外代数 \(\bigwedge (\mathbf{M}^{*})\) （§ 7，第 1 号）。将在 § 11 第 9 号中看到，对所有 \(x\in \mathbf{M}\) ，右内积 \(i(x)\) 是 A 的一个次数为 \(-1\) 的反导子。*

例 5. 回到第 2 号定义 1 的一般情形，设 \(\overline{\mathbf{K}}\) 是另一个交换环，且 \(\rho: \mathbf{K} \to \overline{\mathbf{K}}\) 是一个环同态；设 \(\overline{\mathbf{A}}, \overline{\mathbf{A}'}\) ， \(\overline{\mathbf{A}''}\) ， \(\overline{\mathbf{B}}, \overline{\mathbf{B}'}\) ， \(\overline{\mathbf{B}''}\) 表示分次 \(\overline{\mathbf{K}}\) -模，它们分别由 \(\mathbf{A}, \mathbf{A}', \mathbf{A}''\) ， \(\mathbf{B}, \mathbf{B}'\) ， \(\mathbf{B}''\) 通过将标量环扩张到 \(\overline{\mathbf{K}}\) 而得到（II，§ 11，第 5 号）；我们由 \(\mu\) ， \(\lambda_1\) 和 \(\lambda_2\) 导出 \(\overline{\mathbf{K}}\) -双线性映射

\[
\bar {\mu}: \overline {{{A}}} \times \overline {{{A}}} ^ {\prime} \rightarrow \overline {{{A}}} ^ {\prime \prime}, \quad \bar {\lambda} _ {1}: \bar {B} \times \overline {{{A}}} ^ {\prime} \rightarrow \bar {B} ^ {\prime \prime}, \quad \bar {\lambda} _ {2}: \overline {{{A}}} \times \bar {B} ^ {\prime} \rightarrow \bar {B} ^ {\prime \prime}
\]

作 \(l_{\bar{K}}\) 与 \(\mu\) ， \(\lambda_{1}\) 和 \(\lambda_{2}\) 相应的 K-线性映射的张量积（II, § 5, 第 1 号）。那么，若 d 是一个 \(\varepsilon\) -导子（次数为 \(\delta\) ，从 \((\mathbf{A}, \mathbf{A}', \mathbf{A}^{\prime\prime})\) 到 \((\mathbf{B}, \mathbf{B}', \mathbf{B}^{\prime\prime})\) ），则映射 \(\bar{d} = d \otimes l_{\bar{K}}\) （从 \(\overline{A} \oplus \overline{A}' \oplus \overline{A}''\) 到 \(\overline{B} \oplus \overline{B}' \oplus \overline{B}''\) ）是一个 \(\varepsilon\) -导子（次数为 \(\delta\) ，从 \((\overline{A}, \overline{A}', \overline{A}^{\prime\prime})\) 到 \((\overline{B}, \overline{B}', \overline{B}'' )\) ）。

例 6. 设 A 为类型 Z 的分次 K-代数；通过如下方式定义一个次数为 0 的分次 K-线性映射 \(d: A \to A\) ：对 \(x_{n} \in \mathbf{A}_{n}(n \in \mathbf{Z})\) ，取 \(d(x_{n}) = nx_{n}\) 。该映射是 A 的一个导子，因为对 \(x_{p} \in A_{p}, x_{q} \in A_{q}\) ，

\[
d (x _ {p} x _ {q}) = (p + q) x _ {p} x _ {q} = d (x _ {p}) x _ {q} + x _ {p} d (x _ {q}).
\]

\subsection{导子的复合}

我们在本号中假设第 2 号的情形（I）成立，即 A、A'、A'' 是三个类型为 Δ 的分次 K-模，并且给定一个 K-双线性映射 μ: A × A' → A''，它对应于一个次数为 0 的分次 K-线性映射 A ⊗K A' → A''。A ⊕ A' ⊕ A'' 的分次自同态 f，若满足 f(A) ⊂ A、f(A') ⊂ A' 且 f(A'') ⊂ A''，则构成分次结合代数 EndgrK(A ⊕ A' ⊕ A'')（§ 3，第 1 号，例 2）的一个分次子代数。特别地，两个ε-导子（关于(A, A', A'')）可以复合，但不应认为两个ε-导子的复合仍是ε-导子。

在每一个类型为 \(\Delta\) 的分次代数 B 上， \(\varepsilon\) -括号（当 \(\varepsilon = 1\) 时简称括号）由公式 (10) 对两个齐次元素 \(u, v\) 定义：

\[
[ u, v ] _ {\varepsilon} = u v - \varepsilon_ {\deg u, \deg v} v u (\text { denoted   simply   by } [ u, v ] \text { if } \varepsilon = 1).
\]

将此映射线性扩张，就得到一个从 B × B 到 B 的 K-双线性映射 \((u, v) \mapsto [u, v]_{\varepsilon}\) 。于是，对 B 中齐次的 u 和 v

\[
[ v, u ] _ {\varepsilon} = - \varepsilon_ {\deg u, \deg v} [ u, v ] _ {\varepsilon}.
\]

将此定义应用于分次代数 \(\operatorname{Endgr}_{\mathbf{K}}(\mathbf{A} \oplus \mathbf{A}' \oplus \mathbf{A}'')\) ，两个分次自同态的 \(\varepsilon\) -括号便由此得到定义。

命题 1. 设 \(d_1, d_2\) 是两个 \(\varepsilon\) -导子（属于 \((\mathbf{A}, \mathbf{A}', \mathbf{A} ' )\) ，次数分别为 \(\delta_1, \delta_2\) ）。则 \(\varepsilon\) -括号

\[
[ d _ {1}, d _ {2} ] _ {\varepsilon} = d _ {1} \circ d _ {2} - \varepsilon_ {\delta_ {1}, \delta_ {2}} d _ {2} \circ d _ {1}
\]

是一个 \(\varepsilon\) -导子（次数为 \(\delta_1 + \delta_2\) ）。此外，若 \(d\) 是一个 \(\varepsilon\) -导子（属于 \((\mathbf{A}, \mathbf{A}', \mathbf{A}'')\) ，次数为 \(\delta\) ），且 \(\varepsilon_{\delta, \delta} = -1\) ，则 \(d^2 = d \circ d\) 是一个导子。

设 \(x \in A\) 是次数为 \(\xi\) 的齐次元素；对所有 \(y \in A'\) ，

\[
\begin{array}{r l} d _ {1} (d _ {2} (x y)) & = ((d _ {1} d _ {2}) (x)) y + \varepsilon_ {\delta_ {1}, \delta_ {2} + \xi} (d _ {2} x) (d _ {1} y) \\ & \quad + \varepsilon_ {\delta_ {2}, \xi} (d _ {1} x) (d _ {2} y) + \varepsilon_ {\delta_ {1} + \delta_ {2}, \xi} x ((d _ {1} d _ {2}) (y)) \end{array}\tag{11}
\]

这里利用了第 1 号的公式 (1)。交换 \(d_1\) 和 \(d_2\) 的角色，并再次利用 (1)（第 1 号）化简后，我们得到

\[
\begin{array}{r l} (d _ {1} d _ {2}) (x y) - \varepsilon_ {\delta_ {1}, \delta_ {2}} (d _ {2} d _ {1}) (x y) = & ((d _ {1} d _ {2}) (x)) y - \varepsilon_ {\delta_ {1}, \delta_ {2}} ((d _ {2} d _ {1}) (x)) y \\ & + \varepsilon_ {\delta_ {1} + \delta_ {2}, \xi} x ((d _ {1} d _ {2}) (y)) \\ & - \varepsilon_ {\delta_ {1}, \delta_ {2}} \varepsilon_ {\delta_ {1} + \delta_ {2}, \xi} x ((d _ {2} d _ {1}) (y)) \end{array}
\]

即，记 \(d = [d_1, d_2]_{\varepsilon}\) 和 \(\delta = \delta_1 + \delta_2\) ,

\[
d (x y) = (d x) y + \varepsilon_ {\delta , \xi} x (d y)
\]

这证明了 \(d\) 是一个 \(\varepsilon\) -导子。

另一方面，若在 (11) 中令 \(d_1 = d_2 = d\) ， \(\delta_1 = \delta_2 = \delta\) 以及 \(\varepsilon_{\delta,\delta} = -1\) ，则由 (1)，此时 \(\varepsilon_{\delta,\delta + \xi} = -\varepsilon_{\delta,\xi}\) ，故得到

\[
d ^ {2} (x y) = \left(d ^ {2} x\right) y + \varepsilon_ {2 5, \xi} x \left(d ^ {2} y\right)
\]

又由于 \(\varepsilon_{2\delta, \xi} = 1\) ，可以看出 \(d^2\) 是一个导子。

推论。假设 \(\Delta = \mathbf{Z}\) 。那么：

\begin{enumerate}
\def\labelenumi{(\roman{enumi})}
\item
  一个反导子的平方是一个导子。
\item
  两个导子之间的括号运算仍是一个导子。
\item
  反导子与偶数次导子的括号运算仍为反导子。
\item
  若 \(d_1\) 和 \(d_2\) 是反导子，则 \(d_1d_2 + d_2d_1\) 是一个导子。
\end{enumerate}

在本号开头的假设下，现在考虑一个有限序列 \(\mathbf{D} = (d_i)_{1 \leqslant i \leqslant n}\) ，它由 \((\mathbf{A}, \mathbf{A}', \mathbf{A}'')\) 的两两可交换的导子组成。对每个多项式 \(\mathrm{P}(\mathbf{X}_1, \ldots, \mathbf{X}_n)\) （属于代数 \(\mathbf{K}[\mathbf{X}_1, \ldots, \mathbf{X}_n]\) ），元素 \(\mathrm{P}(d_1, \ldots, d_n)\) （ \(\operatorname{Endgr}_{\mathbf{K}}(\mathbf{A} \oplus \mathbf{A}' \oplus \mathbf{A}'')\) 中的）便有定义（§ 2，第 9 号）；其简写记号为 \(\mathrm{P}(\mathbf{D})\) 。

命题 2. 在上述假设与记号下，考虑 \(2n\) 个未定元 \(\mathrm{T}_1, \ldots, \mathrm{T}_n\) ， \(\mathrm{T}_1', \ldots, \mathrm{T}_n'\) ，并且对每个多项式 \(\mathbf{F} \in \mathbf{K}[\mathbf{X}_1, \ldots, \mathbf{X}_n]\) ，记 \(\mathbf{F}(\mathbf{T}) = \mathbf{F}(\mathbf{T}_1, \ldots, \mathbf{T}_n)\) ， \(\mathbf{F}(\mathbf{T}') = \mathbf{F}(\mathbf{T}_1', \ldots, \mathbf{T}_n')\) 以及

\[
\mathrm{F} (\mathbf {T} + \mathbf {T} ^ {\prime}) = \mathrm{F} (\mathrm{T} _ {1} + \mathrm{T} _ {1} ^ {\prime}, \dots , \mathrm{T} _ {n} + \mathrm{T} _ {n} ^ {\prime}).
\]

假设

\[
\mathbf {P} (\mathbf {T} + \mathbf {T} ^ {\prime}) = \sum_ {i} \mathbf {Q} _ {i} (\mathbf {T}) \mathbf {R} _ {i} (\mathbf {T} ^ {\prime})
\]

其中 \(\mathbf{Q}_i\) 和 \(\mathbf{R}_i\) 属于 \(\mathbf{K}[\mathbf{X}_1,\ldots ,\mathbf{X}_n]\) 。那么，对于所有 \(x\in \mathbf{A}\) 和 \(y\in \mathbf{A}'\) ,

\[
\mathbf {P} (\mathbf {D}) (x y) = \sum_ {i} \left(\mathbf {Q} _ {i} (\mathbf {D}) x\right) (\mathbf {R} _ {i} (\mathbf {D}) y).\tag{12}
\]

我们引入另外 \(n\) 个未定元 \(T_1'', \ldots, T_n''\) ，并考虑多项式代数 \(K[T_1, \ldots, T_n, T_1', \ldots, T_n', T_1'', \ldots, T_n''] = B\) ；另一方面，我们考虑 K-模 \(M\) ，它由 \(A \times A'\) 到 \(A''\) 的双线性映射组成；在 \(M\) 上定义一个 B-模结构如下：对每个 K-双线性映射 \(f \in M\) 以及 \(1 \leqslant i \leqslant n\) ，

\[
\left\{ \begin{array}{l} (\mathrm{T} _ {i} f) (a, a ^ {\prime}) = f (d _ {i} a, a ^ {\prime}) \\ (\mathrm{T} _ {i} ^ {\prime} f) (a, a ^ {\prime}) = f (a, d _ {i} a ^ {\prime}) \\ (\mathrm{T} _ {i} ^ {\prime \prime} f) (a, a ^ {\prime}) = d _ {i} (f (a, a ^ {\prime})). \end{array} \right.\tag{13}
\]

由于诸 \(d_{i}\) 两两可交换，可见对每个多项式 \(\mathbf{F} \in \mathbf{K}[\mathbf{X}_{1}, \ldots, \mathbf{X}_{n}]\) ，有 \((\mathbf{F}(\mathbf{T})f)(a, a') = f(\mathbf{F}(\mathbf{D})a, a')\) ，

\[
(\mathbf {F} (\mathbf {T} ^ {\prime}) f) (a, a ^ {\prime}) = f (a, \mathbf {F} (\mathbf {D}) a ^ {\prime})
\]

以及 \((\mathbf{F}(\mathbf{T}^{\prime \prime})f) = \mathbf{F}(\mathbf{D})(f(a,a^{\prime}))\) 。因此，为证明 (12)，只需证明

\[
(\mathbf {P} (\mathbf {T} ^ {\prime \prime}) - \sum_ {i} \mathbf {Q} _ {i} (\mathbf {T}) \mathbf {R} _ {i} (\mathbf {T} ^ {\prime})) \cdot \mu = 0\tag{14}
\]

或等价地 \((\mathrm{P}(\mathbf{T}^{\prime\prime}) - \mathrm{P}(\mathbf{T} + \mathbf{T}^{\prime}))\) . \(\mu = 0\) 在 B-模 M 中成立。现在，诸 \(d_{i}\) 是导子这一假设也可以表述为：对 \(1 \leqslant i \leqslant n\) ，

\[
\left(\mathrm{T} _ {i} ^ {\prime \prime} - \mathrm{T} _ {i} - \mathrm{T} _ {i} ^ {\prime}\right). \mu = 0\tag{15}
\]

在 B-模 M 中成立。依次考虑多项式

\[
\mathrm{P} \left(\mathrm{T} _ {1} ^ {\prime \prime}, \mathrm{T} _ {2} ^ {\prime \prime}, \dots , \mathrm{T} _ {n} ^ {\prime \prime}\right) - \mathrm{P} \left(\mathrm{T} _ {1} + \mathrm{T} _ {1} ^ {\prime}, \mathrm{T} _ {2} ^ {\prime \prime}, \dots , \mathrm{T} _ {n} ^ {\prime \prime}\right)
\]

\[
\mathrm{P} \left(\mathrm{T} _ {1} + \mathrm{T} _ {1} ^ {\prime}, \mathrm{T} _ {2} ^ {\prime \prime}, \dots , \mathrm{T} _ {n} ^ {\prime \prime}\right) - \mathrm{P} \left(\mathrm{T} _ {1} + \mathrm{T} _ {1} ^ {\prime}, \mathrm{T} _ {2} + \mathrm{T} _ {2} ^ {\prime}, \dots , \mathrm{T} _ {n} ^ {\prime \prime}\right)
\]

\[
\mathrm{P} \left(\mathrm{T} _ {1} + \mathrm{T} _ {1} ^ {\prime}, \dots , \mathrm{T} _ {n - 1} + \mathrm{T} _ {n - 1} ^ {\prime}, \mathrm{T} _ {n} ^ {\prime \prime}\right) - \mathrm{P} \left(\mathrm{T} _ {1} + \mathrm{T} _ {1} ^ {\prime}, \dots , \mathrm{T} _ {n - 1} + \mathrm{T} _ {n - 1} ^ {\prime}, \mathrm{T} _ {n} + \mathrm{T} _ {n} ^ {\prime}\right)
\]

就可以看出，差 \(\mathbf{P}(\mathbf{T}^{\prime \prime}) - \mathbf{P}(\mathbf{T} + \mathbf{T}^{\prime})\) 可以写成如下形式

\[
\sum_ {i = 1} ^ {n} \left(\mathbf {T} _ {i} ^ {\prime \prime} - \mathbf {T} _ {i} - \mathbf {T} _ {i} ^ {\prime}\right) \mathrm{G} _ {i} (\mathbf {T}, \mathbf {T} ^ {\prime}, \mathbf {T} ^ {\prime \prime})
\]

其中 \(G_{i}\) 是 B 的元素。因此关系 (14) 是关系 (15) 的直接推论。

推论（莱布尼茨公式）。设 \(d_{i}\) ( \(1 \leqslant i \leqslant n\) ) 是 \(n\) 个导子（属于 \((\mathbf{A}, \mathbf{A}', \mathbf{A}'' )\) ，两两可交换）。对所有 \(\alpha = (\alpha_{1}, \ldots, \alpha_{n}) \in \mathbf{N}^{n}\) ，我们记

\[
d ^ {\alpha} = d _ {1} ^ {\alpha_ {1}} d _ {2} ^ {\alpha_ {2}} \dots d _ {n} ^ {\alpha_ {n}}.\tag{16}
\]

那么，对 \(x \in \mathbf{A}\) 和 \(y \in \mathbf{A}'\) ，

\[
d ^ {\alpha} (x y) = \sum_ {\beta + \gamma = \alpha} ((\beta , \gamma)) d ^ {\beta} (x) d ^ {\gamma} (y)\tag{17}
\]

其中我们写下了（采用本章开头引入的记号）

\[
((\beta , \gamma)) = (\beta + \gamma)! / (\beta ! \gamma!).\tag{18}
\]

这由多项式公式（I，§ 8，第 2 号）

\[
(\mathbf {T} + \mathbf {T} ^ {\prime}) ^ {\alpha} = \sum_ {\beta + \gamma = \alpha} ((\beta , \gamma)) \mathbf {T} ^ {\beta} \mathbf {T} ^ {\prime \gamma}
\]

以及命题 2 立即得出。

\subsection{代数 A 到 A-模的导子}

我们在本号中假设第 2 号的情形 (II) 成立。那么存在一个分次 K-代数 A 和一个分次 K-模 E，以及两个次数为 0 的 K-线性映射

\[
\mathbf {E} \otimes_ {\mathbf {K}} \mathbf {A} \rightarrow \mathbf {E}, \quad \mathbf {A} \otimes_ {\mathbf {K}} \mathbf {E} \rightarrow \mathbf {E}
\]

记作

\[
x \otimes a \mapsto x. a \quad \text { and } \quad a \otimes x \mapsto a. x \quad \text { for   } a \in A \text {   and   } x \in E.
\]

命题 3. 设 \(d:A\to E\) 是一个 \(\varepsilon\) -导子（次数为 \(\delta\) ）。则 \(\operatorname{Ker}(d)\) 是 A 的一个分次子代数；若 A 有单位元，则它属于 \(\operatorname{Ker}(d)\) 。

显然 \(\operatorname{Ker}(d)\) 是 A 的一个分次子 K-模；此外，第 2 号的关系 (8) 表明，若 x 和 y 是属于 \(\operatorname{Ker}(d)\) 的两个齐次元素，则 \(d(xy) = 0\) ，从而 \(xy \in \operatorname{Ker}(d)\) 。最后，若 A 有单位元 1（次数为 0，参见 § 3 第 1 号），则把第 2 号的关系 (8) 中的 x 和 y 都换成 1，就得到 \(d(1) = d(1) + d(1)\) ，从而 \(d(1) = 0\) 。

推论。设 \(d_1\) 和 \(d_2\) 是两个 \(\varepsilon\) -导子（从 \(A\) 到 \(E\) ，次数同为 \(\delta\) ）。若 \(d_1\) 和 \(d_2\) 在代数 \(A\) 的一个生成系的每个元素上取相同的值，则 \(d_1 = d_2\) 。

\(d_{1} - d_{2}\) 是一个 \(\varepsilon\) -导子（次数为 \(\delta\) ），因此 \(\operatorname{Ker}(d_1 - d_2)\) 是 A 的一个子代数，它包含 A 的一个生成系，从而等于 A。

命题 4. 设 \(d:A\to E\) 是一个 \(\varepsilon\) -导子（次数为 \(\delta\) ）。假设 A 有单位元 1，并设 \(x\) 是 A 的一个齐次元素，它在 A 中有逆元 \(x^{-1}\) 。则

\[
d (x ^ {- 1}) = - \varepsilon_ {\delta , \deg (x)} x ^ {- 1} ((d x) x ^ {- 1}) = - \varepsilon_ {\delta , \deg (x)} (x ^ {- 1} (d x)) x ^ {- 1}.\tag{19}
\]

我们有 \(d(xx^{-1}) = d(1) = 0\) （命题 3），由此得出

\[
(d x) x ^ {- 1} + \varepsilon_ {\delta , \deg (x)} x (d (x ^ {- 1})) = 0
\]

这证明了 (19) 中的第一个公式。另一方面， \(x^{-1}\) 是次数为 \(-\deg(x)\) 的齐次元素，且由第 1 号的公式 (1) 有 \(\varepsilon_{\delta,\deg(x)} = \varepsilon_{\delta,-\deg(x)}\) ；写出 \(d(x^{-1}x) = 0\) ，即类似地得到 (19) 的第二个公式。

命题 5. 设 A 是一个整环，并令 L 为其分式域。A 到 L 上向量空间 E（视为 A-模）的每个导子都可以唯一地扩张为 L 到 E 的导子。

设 \(d\) 是 \(\mathbf{A}\) 到 \(\mathbf{E}\) 的一个导子， \(\bar{d}\) 是 \(\mathbf{L}\) 到 \(\mathbf{E}\) 的一个导子且扩张 \(d\) ；那么，对 \(u \in \mathbf{A}, v \in \mathbf{A}, v \neq 0\) ，必然有，根据 (19)，

\[
\bar {d} (u / v) = v ^ {- 1} d u - u v ^ {- 2} d v\tag{20}
\]

这就证明了 \(\bar{d}\) 的唯一性。反之，我们证明 \(\bar{d}\) 可以由公式 (20) 定义；首先必须验证，若 \(u / v = u' / v'\) ，则当把 \(u\) 换成 \(u'\) 、把 \(v\) 换成 \(v'\) 时，(20) 右端的值不改变。现在， \(uv' = vu'\) ，因此 \(v'(du) + u(dv') = v(du') + u'(dv)\) ，从而 \(v'(du - uv^{-1}dv) = v(du' - u'v'^{-1}dv')\) ，因为 \(uv'v^{-1} = u'\) 且 \(u'v'^{-1}v = u\) 。于是定义了一个映射 \(\bar{d}: L \to E\) ，它扩张 \(d\) ；立即可验证它是 K-线性的且是一个导子。

命题 6. 设 A 是一个含单位元的结合分次 K-代数，E 是一个分次 (A, A)-双模。若 \(d: \mathbf{A} \to \mathbf{E}\) 是一个 \(\varepsilon\) -导子（次数为 \(\delta\) ），则对 A 的齐次元素的每个有限序列 \((x_i)_{1 \leqslant i \leqslant n}\) （次数分别为 \(\xi_i\) ， \(1 \leqslant i \leqslant n\) ），

\[
d \left(x _ {1} x _ {2} \dots x _ {n}\right) = \sum_ {i = 1} ^ {n} \varepsilon_ {\delta, \xi_ {1} + \dots + \xi_ {i - 1}} x _ {1} \dots x _ {i - 1} \left(d x _ {i}\right) x _ {i + 1} \dots x _ {n}.\tag{21}
\]

公式 (21) 对 \(n = 0\) 是平凡的，并对 \(n\) 用归纳法证明，其中用到 (4)（第 2 号）。

推论。设 A 是一个含幺交换结合代数，E 是一个 A-模。若 \(d: \mathbf{A} \to \mathbf{E}\) 是一个导子，则对于每个整数 \(n \geqslant 0\) ,

\[
d (x ^ {n}) = n x ^ {n - 1} (d x) \quad \text { for all } x \in \mathbf {A}.\tag{22}
\]

只需给 A 赋予平凡分次，并在 (21) 中取所有 \(x_{i}\) 等于 \(x\) 即可。

我们回到一个 \(\varepsilon\) -导子 \(d: \mathbf{A} \to \mathbf{E}\) （次数为 \(\delta\) ）的一般情形。设 \(Z_{\varepsilon}\) 是这样的 \(a \in \mathbf{A}\) 组成的集合：使得对每个齐次分量 \(a_{\alpha}\) （属于 \(a\) ，次数为 \(\alpha\) ），对每个齐次的 \(x\) （在 \(\mathbf{E}\) 中），

\[
x a _ {\alpha} = \varepsilon_ {\alpha , \deg (x)} a _ {\alpha} x.\tag{23}
\]

若 A 是一个含单位元的结合分次代数，且 E 是一个分次 (A, A)-双模，则由该定义立即得出 \(Z_{\varepsilon}\) 是A的一个分次子代数，且包含单位元。

命题 7. 设 A 是一个含单位元的结合分次代数，E 是一个分次 (A, A)-双模。令 \(d: \mathbf{A} \to \mathbf{E}\) 是一个 \(\varepsilon\) -导子（次数为 \(\delta\) ），并令 \(a\) 是 \(Z_{\varepsilon}\) 的一个次数为 \(\alpha\) 的齐次元素。则映射 \(x \mapsto a(dx)\) 是一个 \(\varepsilon\) -导子（次数为 \(\delta + \alpha\) ）。

我们记 \(d'(x) = a(dx)\) ；对齐次元素 \(x\) （次数为 \(\xi\) ， \(A\) 中的）以及 \(y \in A\) ，根据 (23) 和 (1)（第 1 号），

\[
\begin{array}{r l} d ^ {\prime} (x y) & = a ((d x) y) + \varepsilon_ {\delta , \xi} a (x (d y)) = (a (d x)) y + \varepsilon_ {\delta + \alpha , \xi} (x a) (d y) \\ & = (d ^ {\prime} x) y + \varepsilon_ {\delta + \alpha , \xi} x (d ^ {\prime} y). \end{array}
\]

命题 7 表明，A 到 E 的 \(\varepsilon\) -导子组成的 K-模是一个分次 \(Z_{\varepsilon}\) -模（类型为 \(\Delta\) ）。

\subsection{代数的导子}

设 A 为分次 K-代数；对每个齐次元素 \(a \in \mathbf{A}\) ，用 \(\operatorname{ad}_{\varepsilon}(a)\) （若不致混淆，简记为 \(\operatorname{ad}(a)\) ）表示 A 到 A 的 K-线性映射

\[
x \mapsto [ a, x ] _ {\varepsilon}
\]

（第 4 号，公式 (10)），它是次数为 deg a 的分次映射。

命题 8. 设 A 为分次 K-代数。

\begin{enumerate}
\def\labelenumi{(\roman{enumi})}
\tightlist
\item
  对每个 \(\varepsilon\) -导子 \(d: \mathbf{A} \to \mathbf{A}\) 以及每个齐次元素 \(a\) （ \(\mathbf{A}\) 中的），
\end{enumerate}

\[
[ d, \mathrm{ad} _ {\varepsilon} (a) ] _ {\varepsilon} = \mathrm{ad} _ {\varepsilon} (d a).\tag{24}
\]

\begin{enumerate}
\def\labelenumi{(\roman{enumi})}
\setcounter{enumi}{1}
\item
  若代数 \(\mathbf{A}\) 是结合的，则 \(\operatorname{ad}_{\varepsilon}(a)\) 是一个 \(\varepsilon\) -导子（ \(\mathbf{A}\) 的，次数为 \(\deg(a)\) ）。
\item
  假设 \(d\) 的次数为 \(\delta\) ，设 \(\alpha = \deg a\) ，并记 \(f = [d, \operatorname{ad}_{\varepsilon}(a)]_{\varepsilon}\) 。对每个齐次元素 \(x \in A\) （次数为 \(\xi\) ），由 (1)（第 1 号），我们有
\end{enumerate}

\[
\begin{array}{r l} f (x) & = d (a x - \varepsilon (\alpha , \xi) x a) - \varepsilon_ {\delta , \alpha} (a (d x) - \varepsilon_ {\alpha , \delta + \xi} (d x) a) \\ & = (d a) x + \varepsilon_ {\delta , \alpha} a (d x) - \varepsilon_ {\alpha , \xi} (d x) a - \varepsilon_ {\delta + \alpha , \xi} x (d a) \\ & \quad - \varepsilon_ {\delta , \alpha} a (d x) + \varepsilon_ {\alpha , \xi} (d x) a \\ & = (d a) x - \varepsilon_ {\delta + \alpha , \xi} x (d a) = [ d a, x ] _ {\varepsilon}. \end{array}
\]

\begin{enumerate}
\def\labelenumi{(\roman{enumi})}
\setcounter{enumi}{1}
\tightlist
\item
  对所有齐次的 \(x\) （次数为 \(\xi\) ）以及所有齐次的 \(y\) （次数为 \(\eta\) ，在 \(\mathbf{A}\) 中），
\end{enumerate}

\[
\begin{array}{r l} \mathrm{ad} _ {\varepsilon} (a) (x y) & = a (x y) - \varepsilon_ {\alpha , \xi + \eta} (x y) a \\ & = (a x - \varepsilon_ {\alpha , \xi} x a) y + \varepsilon_ {\alpha , \xi} x (a y - \varepsilon_ {\alpha , \eta} y a) \\ & = \mathrm{ad} _ {\varepsilon} (a) (x). y + \varepsilon_ {\alpha , \xi} x. \mathrm{ad} _ {\varepsilon} (a) (y) \end{array}
\]

这里用到了 (1) 以及 A 的结合性。

当 A 是结合的时， \(\operatorname{ad}_{\varepsilon}(a)\) 被称为由 a 定义的 A 的内 \(\varepsilon\) -导子。

推论。设 A 是一个结合的分次代数。对于 A 的两个齐次元素 a, b，

\[
[ \mathrm{ad} _ {\varepsilon} (a), \mathrm{ad} _ {\varepsilon} (b) ] _ {\varepsilon} = \mathrm{ad} _ {\varepsilon} ([ a, b ] _ {\varepsilon}).\tag{25}
\]

只需在 (24) 中把 \(d\) 换成 \(\operatorname{ad}_{\varepsilon}(a)\) ，并把 \(\operatorname{ad}_{\varepsilon}(a)\) 换成 \(\operatorname{ad}_{\varepsilon}(b)\) 。

若 \(\deg a = \alpha\) , \(\deg b = \beta\) ，则公式 (25) 等价于下述关系：对每个齐次元素 \(c \in \mathbf{A}\) （次数为 \(\gamma\) ），

\[
\varepsilon_ {\alpha , \gamma} [ a, [ b, c ] _ {\varepsilon} ] _ {\varepsilon} + \varepsilon_ {\beta , \alpha} [ b, [ c, a ] _ {\varepsilon} ] _ {\varepsilon} + \varepsilon_ {\gamma , \beta} [ c, [ a, b ] _ {\varepsilon} ] _ {\varepsilon} = 0\tag{26}
\]

称为雅可比恒等式。

\subsection{函子性质}

在本号中，所有代数均假设为结合的且含单位元的，并且每个代数同态均假设为含单位元的。

命题 9. 设 A、B 为两个分次 K-代数，E 为 (A, A)-双模，F 为分次 (B, B)-双模；设 \(\rho: \mathbf{A} \to \mathbf{B}\) 是一个分次代数同态，并且 \(\theta: \mathbf{E} \to \mathbf{F}\) 是 A-双模（相对于 \(\rho\) ）的一个分次 A-同态，次数为 0。则：

\begin{enumerate}
\def\labelenumi{(\roman{enumi})}
\item
  对每个 \(\varepsilon\) -导子 \(d':\mathbf{B}\to\mathbf{F}\) ， \(d'\circ\rho:\mathbf{A}\to\rho^{*}(\mathbf{F})\) 是一个同次数的 \(\varepsilon\) -导子。
\item
  对每个 \(\varepsilon\) -导子 \(d: \mathbf{A} \to \mathbf{E}\) ， \(\theta \circ d: \mathbf{A} \to \rho^{*}(\mathbf{F})\) 是一个同次数的 \(\varepsilon\) -导子。
\end{enumerate}

这两个断言直接从这些关系中得出。

\[
\begin{array}{l} d ^ {\prime} (\rho (x y)) = d ^ {\prime} (\rho (x) \rho (y)) = d ^ {\prime} (\rho (x)) \rho (y) + \varepsilon_ {\delta ^ {\prime}, \xi} \rho (x) d ^ {\prime} (\rho (y)) \\ \theta (d (x y)) = \theta ((d x) y + \varepsilon_ {\delta, \xi} x (d y)) = \theta (d x) \rho (y) + \varepsilon_ {\delta, \xi} \rho (x) \theta (d y) \end{array}
\]

其中 \(x \in \mathbf{A}\) 是次数为 \(\xi\) 的齐次元素， \(y \in \mathbf{A}\) ，而 \(\delta\) 与 \(\delta'\) 分别表示 \(d\) 与 \(d'\) 的次数。

推论。设 S 是代数 A 的一个生成系。为使 \(d' \circ \rho = \theta \circ d\) ，必要且充分条件是 \(d'(\rho(x)) = \theta(d(x))\) 对所有 \(x \in S\) 成立。

这是命题 9 以及第 5 号命题 3 的推论的直接推论。

在命题 9 的条件下，我们知道 B（借助 \(\rho\) ）具有一个 (A, A)-双模结构（II，§ 1，第 14 号，例 1）。

命题 10. 在命题 9 的条件下，要使 \(\varepsilon\) -导子 \(d': \mathbf{B} \to \mathbf{F}\) 对 B 与 \(\rho_*(\mathbf{F})\) 上的左（相应地，右）A-模结构是 A-线性的，必要且充分条件是 \(d'\) 在 B 的子代数 \(\rho(\mathbf{A})\) 上为零。

我们对左 A-模结构进行证明。对 \(a \in \mathbf{A}\) ， \(b \in \mathbf{B}\) ，

\[
d ^ {\prime} (\rho (a) b) = d ^ {\prime} (\rho (a)) b + \rho (a) d ^ {\prime} b
\]

于是，若 \(d' \circ \rho = 0\) ，则 \(d'\) 对 B 与 \(\rho^*(\mathbf{F})\) 上的左 A-模结构是线性的。反之，若确实如此，则特别地

\[
d ^ {\prime} (\rho (a)) = d ^ {\prime} (\rho (a). 1) = \rho (a) d ^ {\prime} (1) = 0
\]

（第 5 号，命题 3）。

特别地，设 \(D_{K}(B,F)\) 表示 B 到 F 的导子组成的 K-模（第 2 号）；其中 A-线性的那些导子，换句话说在 \(\rho(A)\) 上为零的那些导子，构成 \(D_{K}(B,F)\) 的一个子 K-模，记作 \(D_{A,\rho}(B,F)\) 或简记作 \(D_{A}(B,F)\) （显然 \(D_{K}(B,F)=D_{K,\phi}(B,F)\) ，其中 \(\phi:K\to B\) 是定义 B 上 K-代数结构的同态）。

现在设 A、B、C 为三个分次 K-代数， \(\rho:A\to B\) ， \(\sigma:B\to C\) 为两个分次代数同态，G 为一个分次 (C, C)-双模；如果 \(D_{A}(B,G)\) , \(D_{B}(C,G)\) 和 \(D_{A}(C,G)\) 表示相应的 K-模 \(D_{A,\rho}(B,\sigma_{*}(G))\) , \(D_{B,\sigma}(C,G)\) 和 \(D_{A,\sigma_{o}\rho}(C,G)\) ，则 \(D_{B}(C,G)\) 显然是 \(D_{A}(C,G)\) 的一个子 K-模，因为 \(\sigma(\rho(A))\subset\sigma(B)\) 。

命题 11. 在上述条件下，存在 K-同态的一个正合序列

\[
0 \rightarrow \mathrm{D} _ {\mathrm{B}} (\mathrm{C}, \mathrm{G}) \xrightarrow {u} \mathrm{D} _ {\mathrm{A}} (\mathrm{C}, \mathrm{G}) \xrightarrow {v} \mathrm{D} _ {\mathrm{A}} (\mathrm{B}, \mathrm{G})\tag{27}
\]

其中 u 是典范单射，v 是同态 \(d \mapsto d \circ \sigma\) （命题 9）。v 的核是这样的导子 \(d: C \to G\) 组成的集合：\(d(\sigma(b)) = 0\) 对所有 \(b \in B\) 成立；它正是 u 的像。

\subsection{导子与代数同态之间的关系}

我们在本号中再次假设第 2 号的情形 (II) 成立，且不假设分次 K-代数 A 是结合的。给定一个元素 \(\delta \in \Delta\) ，考虑分次 K-模 E( \(\delta\) ) （II，§ 11，第 2 号），它使得

\[
(\mathbf {E} (\delta)) _ {\mu} = \mathbf {E} _ {\mu + \delta}
\]

对所有 \(\mu \in \Delta\) 成立。我们在分次 K-模 \(\mathbf{A} \oplus \mathbf{E}(\delta)\) 上定义一个分次 K-代数结构：对每个齐次元素 \(a \in \mathbf{A}\) 以及任意元素 \(a' \in \mathbf{A}, x, x'\) （ \(\mathbf{E}(\delta)\) 中的），规定

\[
(a, x) \left(a ^ {\prime}, x ^ {\prime}\right) = \left(a a ^ {\prime}, x. a ^ {\prime} + \varepsilon_ {\delta, \deg (a)} a. x ^ {\prime}\right);\tag{28}
\]

验证这个乘法给出一个分次环结构是直接的。

投影 \(p:(a,x)\mapsto a\) 称为代数 \(\mathbf{A}\oplus\mathbf{E}(\delta)\) 的增广（映射），并且是一个分次代数同态。分次 K-线性映射 \(g:\mathbf{A}\to\mathbf{A}\oplus\mathbf{E}(\delta)\) （次数为 0 ）使得复合

\[
\mathrm{A} \xrightarrow {g} \mathrm{A} \oplus \mathrm{E} (\delta) \xrightarrow {p} \mathrm{A}
\]

是恒同映射 \(l_{A}\) 的那些映射，正是形如 \(x \mapsto (x, f(x))\) 的映射，其中 \(f: A \to E\) 是次数为 \(\delta\) 的分次 K-线性映射。

命题 12. 为使分次 K-线性映射 \(f:A\to E\) （次数为 \(\delta\) ）是一个 \(\varepsilon\) -导子，必要且充分条件是：映射 \(x\mapsto(x,f(x))\) （从 A 到 \(\mathbf{A}\oplus\mathbf{E}(\delta)\) ）是一个分次 K-代数同态。

利用如下事实：对齐次的 \(x\) （ \(\mathbf{A}\) 中的）以及 \(y \in \mathbf{A}\) ，

\[
(x y, f (x y)) = (x, f (x)). (y, f (y)),
\]

并考虑到(28)，我们得到等价关系

\[
f (x y) = f (x). y + \varepsilon_ {\delta , \deg (x)} x. f (y),
\]

由此得出该命题。

命题 13. 为使代数 \(\mathbf{A} \oplus \mathbf{E}(\delta)\) 是结合的且有单位元的，必要且充分的是： \(\mathbf{A}\) 是结合的且有单位元的，并且映射 \((a, x) \mapsto a.x\) 和 \((a, x) \mapsto x\) . \(a\) 在 \(\mathbf{E}\) 上定义一个 \((\mathbf{A}, \mathbf{A})\) -双模结构； \(\mathbf{A} \oplus \mathbf{E}(\delta)\) 的单位元则为 \((1, 0)\) 。

若把元素 \((u,m)\in\mathbf{A}\oplus\mathbf{E}(\delta)\) 写成该代数的单位元，则立即发现 u 必为 A 的单位元；写出 \((u,m).(0,x)=(0,x).(u,m)=(0,x)\) ，我们得到 u.x=x.u=x （对 \(x\in E\) ），并且写出 \((u,m).(u,0)=(u,0).(u,m)=(u,0)\) ，我们得到 m=0。当 \(\mathbf{A}\oplus\mathbf{E}(\delta)\) 结合时，A 的结合性由增广是一个满同态这一事实得出。条件 \((x.a') .a''=x.(a'.a'')\) 等价于 \(((0,x)(a',0))(a'',0)=(0,x)((a',0)(a'',0))\) ，类似地，条件 \(a.(a'.x)=(a.a').x\) 等价于

\[
(a, 0) ((a ^ {\prime}, 0) (0, x)) = ((a, 0) (a ^ {\prime}, 0)) (0, x);
\]

最后，条件 \(a \cdot (x \cdot a') = (a \cdot x) \cdot a'\) 等价于

\[
(a, 0) ((0, x) (a ^ {\prime}, 0)) = ((a, 0) (0, x)) (a ^ {\prime}, 0).
\]

\subsection{导子的扩张}

命题 14. 设 A 为交换环，M 为 A-模，B 为 A-代数 \(\mathsf{T}(\mathbf{M})\) （相应地 \(\mathsf{S}(\mathbf{M})\) ，相应地 \(\bigwedge (\mathbf{M}))\) ），且 E 为一个 (B, B)-双模。设 \(d_0: \mathbf{A} \to \mathbf{E}\) 是环 A 到 A-模 E 中的一个导子，并设 \(d_{1}:M\to E\) 是一个加法群同态，使得对所有 \(a\in A\) 以及所有 \(x\in M\) ，

\[
d _ {1} (a x) = a d _ {1} (x) + d _ {0} (a). x,\tag{29}
\]

并且进一步，当 \(\mathbf{B} = \mathbf{S}(\mathbf{M})\) 时，

\[
x. d _ {1} (y) + d _ {1} (x). y = y. d _ {1} (x) + d _ {1} (y). x\tag{30}
\]

对所有 \(x, y\) （ \(\mathbf{M}\) 中的）成立，并且，当 \(\mathbf{B} = \bigwedge (\mathbf{M})\) 时，

\[
x. d _ {1} (x) + d _ {1} (x). x = 0\tag{31}
\]

对所有 \(x \in \mathbf{M}\) 成立。那么存在且仅存在一个导子 \(d\) （从 \(\mathbf{B}\) ——视为一个 \(\mathbf{Z}\) -代数——到 \((\mathbf{B}, \mathbf{B})\) -双模 \(\mathbf{E}\) ），使得 \(d|_{\mathbf{A}} = d_0\) 且 \(d|M = d_1\) 。

我们在 \(\mathbf{Z}\) -模 \(\mathbf{B} \oplus \mathbf{E}\) 上取由下式定义的结合 \(\mathbf{Z}\) -代数结构

\[
(b, t) (b ^ {\prime}, t ^ {\prime}) = (b b ^ {\prime}, b t ^ {\prime} + t b ^ {\prime})
\]

它以 (1, 0) 为单位元（第 8 号，命题 13）。在典范单射 \(t \mapsto (0, t)\) 下，E 等同于 B ⊕ E 的一个双边理想，使得 \(\mathrm{E}^2 = \{0\}\) 。另一方面，映射 \(h_0 \colon \mathbf{B} \oplus \mathbf{E}\) （由 \(h_0(a) = (a, d_0(a))\) 定义）是一个有单位元的环同态（第 8 号，命题 12）；在此映射下，B ⊕ E 于是成为一个 A-代数。此外，若对所有 \(x \in \mathbf{M}\) 记 \(h_1(x) = (x, d_1(x))\) ，则由 \(h_0\) 的定义以及 (29) 可知 \(h_1(ax) = h_0(a)h_1(x)\) ；换言之， \(h_1\) 是从 M 到 B ⊕ E 的一个 A-线性映射。于是存在且仅存在一个 A-代数同态 \(h \colon \mathbf{B} \to \mathbf{B} \oplus \mathbf{E}\) ，使得 \(h \mid \mathbf{M} = h_1\) （且必然有 \(h \mid \mathbf{A} = h_0\) ）：因为，若 \(\mathbf{B} = \mathsf{T}(\mathbf{M})\) ，这由 § 5 第 1 号的命题 1 得出；若 \(\mathbf{B} = \mathbb{S}(\mathbf{M})\) ，条件 (30) 表明 \(h(x)h(y) = h(y)h(x)\) 对 M 中所有 \(x, y\) 成立，结论由 § 6 第 1 号的命题 2 得出；最后若 \(\mathbf{B} = \bigwedge (\mathbf{M})\) ，条件 (31) 表明 \((h(x))^2 = 0\) 对所有 \(x \in \mathbf{M}\) 成立，又因 \(x \wedge x = 0\) ，结论由 § 7 第 1 号的命题 1 得出。同态 \(h\) 使得 \(p \circ h \colon \mathbf{B} \to \mathbf{B}\) 与增广 \(\rho \colon \mathbf{B} \oplus \mathbf{E} \to \mathbf{B}\) 的复合是恒等同态 \(l_B\) ，因为 \(p \circ h\) 与 \(l_B\) 在 A 的元素和 M 的元素上按定义一致，而这些元素组成的集合是 B 的一个生成系。因此我们可以写出 \(h(b) = (b, d(b))\) （对所有 \(b \in B\) ），且映射 \(b \mapsto d(b)\) （从 B 到 E 的）依第 8 号的命题 12 是一个具有所需性质的导子。

推论. 设 \(\mathbf{M}\) 是一个分次 \(\mathbf{K}\) -模（类型为 \(\Delta\) ）； \(\mathbf{K}\) -代数 \(\mathsf{T}(\mathbf{M})\) ， \(\mathsf{S}(\mathbf{M})\) 和 \(\bigwedge (\mathbf{M})\) 被赋予相应的类型为 \(\Delta' = \Delta \times \mathbf{Z}\) 的分次（ § 5，第 5 号，命题 7 ； § 6，第 6 号，命题 10 及 § 7，第 7 号，命题 11）。另一方面， \(\mathbf{M}\) 被赋予类型为 \(\Delta'\) 的分次，使得 \(\mathbf{M}_{\alpha,1} = \mathbf{M}_{\alpha}\) 对所有 \(\alpha \in \Delta\) 成立，且 \(\mathbf{M}_{\alpha,n} = \{0\}\) 对 \(\alpha \in \Delta\) 和 \(n \neq 1\) 成立。设 \(\varepsilon'\) 是 \(\Delta'\) 上的一个交换因子。

\begin{enumerate}
\def\labelenumi{(\roman{enumi})}
\item
  设 \(\mathbf{E}\) 是一个分次的（左和右） \(\mathsf{T}(\mathbf{M})\) -双模（类型为 \(\Delta'\) ）；对所有 \(\delta \in \Delta\) 以及每个整数 \(n \in \mathbf{Z}\) ，每个分次 \(\mathbf{K}\) -线性映射 \(f: \mathbf{M} \to \mathbf{E}\) （次数为 \(\delta_1' = (\delta, n)\) ）都可以唯一地扩张为一个 \(\varepsilon'\) -导子 \(d: \mathsf{T}(\mathbf{M}) \to \mathbf{E}\) （次数为 \(\delta'\) ）。
\item
  设 \(\mathbf{E}\) 是一个分次 \(\mathbb{S}(\mathbf{M})\) -模（类型为 \(\Delta'\) ）；为使分次 \(\mathbf{K}\) -线性映射 \(f: \mathbf{M} \to \mathbf{E}\) （次数为 \(\delta'\) ）可以扩张为一个 \(\varepsilon'\) -导子 \(d: \mathbb{S}(\mathbf{M}) \to \mathbf{E}\) （次数为 \(\delta'\) ），必要且充分条件是：对每个有序对 \((x, y)\) （由 \(\mathbf{M}\) 的齐次元素组成），
\end{enumerate}

\[
x. f (y) + \varepsilon_ {\delta^ {\prime}, (\deg (y), 1)} ^ {\prime} y. f (x) = y. f (x) + \varepsilon_ {\delta^ {\prime}, (\deg (x), 1)} ^ {\prime} x. f (y).\tag{33}
\]

该 \(\varepsilon'\) -导子 \(d\) 则是唯一的。

\begin{enumerate}
\def\labelenumi{(\roman{enumi})}
\setcounter{enumi}{2}
\tightlist
\item
  设 \(\mathbf{E}\) 是一个分次的（左和右） \(\bigwedge (\mathbf{M})\) -双模（类型为 \(\Delta'\) ）；为使分次 K-线性映射 \(f: \mathbf{M} \to \mathbf{E}\) （次数为 \(\delta'\) ）可以扩张为一个 \(\varepsilon'\) -导子
\end{enumerate}

\[
d: \bigwedge (\mathbf {M}) \rightarrow \mathbf {E}
\]

（次数为 \(\delta'\) ），必要且充分条件是：对每个齐次元素 \(x\) （属于 \(\mathbf{M}\) ），

\[
x. f (x) + \varepsilon_ {\delta ^ {\prime}, (\deg (x), 1)} ^ {\prime} f (x). x = 0.\tag{34}
\]

该 \(\varepsilon'\) -导子 \(d\) 则是唯一的。

第 2 号的注记 2 应用于如下情形：E 上的外部 B-模律之一（其中 B 等于 T(M)、S(M) 或 ∧(M)）被修改；如此修改后的外部律依 (1) （第 1 号）仍是一个 B-模律，且由此在 E 上得到的 B-模律与另一个 B-模结构仍然相容。这时只需将命题 14 应用于 A = K 与 \(d_{0} = 0\) 的情形。

例 (1). 在命题 14 的应用中，注意若 \(d_0 = 0\) ，条件 (29) 仅表示 \(d_1\) 是 A-线性的。若特别取 \(\mathbf{E} = \mathbf{B}\) 以及由 B 上的环结构导出的 (B, B)-双模结构，则当 \(d_1\) 取为某个自同态 \(s\) （ \(\mathbf{M}\) 的）与典范单射 \(\mathbf{M} \to \mathbf{B}\) 的复合时，条件 (30) 和 (31) 自动满足；这对 (30) 是显然的，因为 \(\mathbb{S}(\mathbf{M})\) 是交换的，而对 (31)，这由以下事实得出： \(x\) 和 \(s(x)\) 在 \(\bigwedge (\mathbf{M})\) 中都是次数 1 的。由此可以看出，每个自同态 \(s\) （ \(\mathbf{M}\) 的）都可以唯一地延拓为一个导子 \(\mathbf{D}_s\) （ \(\mathbb{T}(\mathbf{M})\) 的，相应地 \(\mathbb{S}(\mathbf{M})\) 的，相应地 \(\bigwedge (\mathbf{M})\) 的），其次数为 0。此外，对于两个自同态 \(s, t\) （ \(\mathbf{M}\) 的），

\[
[ \mathbf {D} _ {s}, \mathbf {D} _ {t} ] = \mathbf {D} _ {[ s, t ]}\tag{35}
\]

因为两边都是 \(\mathsf{T}(\mathbf{M})\) （相应地 \(\mathbb{S}(\mathbf{M})\) ，相应地 \(\bigwedge (\mathbf{M})\) ）的导子，且都等于 \([s, t]\) （在 \(\mathbf{M}\) 上）。

\(D_{s}\) 的表达式由第 5 号的公式 (21) 得到，该公式对 M 中的 \(x_{1}, x_{2}, \ldots, x_{n}\) 分别给出：

\[
\left\{ \begin{array}{l} \mathrm{D} _ {s} (x _ {1} \otimes x _ {2} \otimes \dots \otimes x _ {n}) \\ \qquad = \sum_ {i = 1} ^ {n} x _ {1} \otimes \dots \otimes x _ {i - 1} \otimes s (x _ {i}) \otimes x _ {i + 1} \otimes \dots \otimes x _ {n} \\ \mathrm{D} _ {s} (x _ {1} x _ {2} \dots x _ {n}) = \sum_ {i = 1} ^ {n} x _ {1} \dots x _ {i - 1} s (x _ {i}) x _ {i + 1} \dots x _ {n} \\ \mathrm{D} _ {s} (x _ {1} \wedge x _ {2} \wedge \dots \wedge x _ {n}) \\ \qquad = \sum_ {i = 1} ^ {n} x _ {1} \wedge \dots \wedge x _ {i - 1} \wedge s (x _ {i}) \wedge x _ {i + 1} \wedge \dots \wedge x _ {n}. \end{array} \right.\tag{36}
\]

在代数 \(\bigwedge (\mathbf{M})\) 的情形，对 \(\mathbf{D}_s\) 有如下的诠释：

命题 15. 若 M 是有限秩 n 的自由 K-模，则对 M 的每个自同态 s ，限制在 \(\bigwedge^{n}(\mathbf{M})\) 上时，导子 \(\mathbf{D}_{s}\) 是比率为 \(\operatorname{Tr}(s)\) 的位似。

设 \((e_j)_{1\leqslant j\leqslant n}\) 是 M 的一个基，并记 \(e = e_1\wedge e_2\wedge \dots \wedge e_n\) 。若

\[
s (e _ {j}) = \sum_ {k = 1} ^ {n} \alpha_ {j k} e _ {k},
\]

(36) 中的第三个公式给出

\[
\mathrm{D} _ {s} (e) = \sum_ {i = 1} ^ {n} e _ {1} \wedge \dots \wedge e _ {i - 1} \wedge s \left(e _ {i}\right) \wedge e _ {i + 1} \wedge \dots \wedge e _ {n} = \left(\sum_ {j = 1} ^ {n} \alpha_ {j j}\right) e.
\]

例 (2). 在命题 14 的推论的第 (iii) 部分中，令 \(\Delta = \{0\}\) ，此时 \(\bigwedge (\mathbf{M})\) 上的分次是通常的类型为 \(\mathbf{Z}\) 的分次；另一方面取 \(\varepsilon(p, q) = (-1)^{pq}\) 。那么，对每个线性形式 \(x^{*} \in \mathbf{M}^{*}\) （ \(\mathbf{M}\) 上的）， \(x \mapsto \langle x, x^{*} \rangle\) 是一个次数为 \(-1\) 的分次 K-线性映射（从 \(\mathbf{M}\) 到 \(\bigwedge (\mathbf{M})\) ），满足关系式 (34)；于是存在一个反导子 \(i(x^{*})\) （ \(\bigwedge (\mathbf{M})\) 上的，次数为 \(-1\) ），使得（依第 5 号的公式 (21)）

\[
i \left(x ^ {*}\right) \left(x _ {1} \wedge \dots \wedge x _ {n}\right)
\]

\[
= \sum_ {i = 1} ^ {n} (- 1) ^ {i - 1} \left\langle x _ {i}, x ^ {*} \right\rangle x _ {1} \wedge \dots \wedge x _ {i - 1} \wedge x _ {i + 1} \wedge \dots \wedge x _ {n}
\]

并且这是将在 § 11 第 9 号公式 (68) 中定义的内积的一个特例。

命题 16. 设 A 为交换 K-代数， \(\mathbf{M}_i\) ( \(1 \leqslant i \leqslant n\) ) 和 P 为 A-模，H 为从 \(\mathbf{M}_1 \times \mathbf{M}_2 \times \cdots \times \mathbf{M}_n\) 到 P 中的 A-多重线性映射所构成的 A-模。假设给定代数 A 的一个 K-导子 \(d_{0}:A\to A\) ，对每个 i 给定一个 K-线性映射 \(d_{i}:M_{i}\to M_{i}\) 以及一个 K-线性映射 D:P→P，使得对 \(1\leqslant i\leqslant n\) ， \((d_{0},d_{i},d_{i})\) 是 \((\mathbf{A},\mathbf{M}_{i},\mathbf{M}_{i})\) 到自身的一个 K-导子，且 \((d_{0},\mathbf{D},\mathbf{D})\) 是 \((\mathbf{A},\mathbf{P},\mathbf{P})\) 到自身的一个 K-导子。则存在一个 K-线性映射 \(D^{\prime}:H\to H\) ，使得 \((d_{0},D^{\prime},D^{\prime})\) 是 \((\mathbf{A},\mathbf{H},\mathbf{H})\) 到自身的一个 K-导子，且

\begin{enumerate}
\def\labelenumi{(\arabic{enumi})}
\setcounter{enumi}{36}
\tightlist
\item
  \(\mathbf{D}(f(x_1, \ldots, x_n))\)
\end{enumerate}

对所有 \(x_{i} \in \mathbf{M}_{i}\) （ \(1 \leqslant i \leqslant n\) ）和 \(f \in \mathbf{H}\) 成立。

我们证明，对 \(f \in H\) ，映射 \(D'f\) （从 \(M_{1} \times M_{2} \times \cdots \times M_{n}\) 到 P 中的，由 (37) 定义）是 A-多重线性的。对 \(a \in A\) ，

\[
\begin{array}{r l} (\mathrm{D} ^ {\prime} f) (a x _ {1}, x _ {2}, \ldots , x _ {n}) & = \mathrm{D} (a f (x _ {1}, \ldots , x _ {n})) - f (d _ {1} (a x _ {1}), x _ {2}, \ldots , x _ {n}) \\ & \quad - a \sum_ {i = 2} ^ {n} f (x _ {1}, \ldots , x _ {i - 1}, d _ {i} x _ {i}, x _ {i + 1}, \ldots , x _ {n}) \end{array}
\]

且根据假设

\[
\mathrm{D} (a f (x _ {1}, \dots , x _ {n})) = (d _ {0} a) f (x _ {1}, \dots , x _ {n}) + a \mathrm{D} (f (x _ {1}, \dots , x _ {n}))
\]

和 \(d_{1}(ax_{1}) = (d_{0}a)x_{1} + a.d_{1}x_{1}\) ，这给出

\[
(\mathrm{D} ^ {\prime} f) (a x _ {1}, x _ {2}, \dots , x _ {n}) = a. (\mathrm{D} ^ {\prime} f) (x _ {1}, \dots , x _ {n})
\]

而关于每个 \(x_{i}\) 的线性性可类似地证明。另一方面，

\[
\begin{array}{r l} (\mathrm{D} ^ {\prime} (a f)) (x _ {1}, \dots , x _ {n}) & = \mathrm{D} (a f (x _ {1}, \dots , x _ {n})) \\ & - \sum_ {i = 1} ^ {n} a f (x _ {1}, \dots , x _ {i - 1}, d _ {i} x _ {i}, x _ {i + 1}, \dots , x _ {n}) \\ & = (d _ {0} a) f (x _ {1}, \dots , x _ {n})) + a \mathrm{D} (f (x _ {1}, \dots , x _ {n})) \\ & - \sum_ {i = 1} ^ {n} a f (x _ {1}, \dots , x _ {i - 1}, d _ {i} x _ {i}, x _ {i + 1}, \dots , x _ {n}) \\ & = (d _ {0} a) f (x _ {1}, \dots , x _ {m}) + a (\mathrm{D} ^ {\prime} f) (x _ {1}, \dots , x _ {m}) \end{array}
\]

换句话说

\[
\mathbf {D} ^ {\prime} (a f) = \left(d _ {0} a\right) f + a (\mathbf {D} ^ {\prime} f)
\]

这便确立了该命题。

例. (3) 将命题 16 应用于情形 \(n = 1\) ， \(\mathbf{M}_1 = \mathbf{M}\) ， \(P = A\) ，则 \(H = M^*\) ，即 \(M\) 的对偶，并且可以看出，由 K-导子 \((d_0, d, d)\) （ \((A, M, M)\) 的）可导出一个 K-导子 \((d_0, d^*, d^*)\) （ \((A, M^*, M^*)\) 的），使得

\[
d _ {0} \langle m, m ^ {*} \rangle = \langle d m, m ^ {*} \rangle + \langle m, d ^ {*} m ^ {*} \rangle\tag{38}
\]

对 \(m \in \mathbf{M}\) 和 \(m^* \in \mathbf{M}^*\) 成立。从 \(\mathbf{M} \oplus \mathbf{M}^*\) 到自身的、等于 \(d\) （在 \(\mathbf{M}\) 上）且等于 \(d^{*}\) （在 \(\mathbf{M}^{*}\) 上）的那个 K-线性映射于是满足条件 (29)，因此存在一个 K-导子 \(\mathbf{D}\) （A-代数 \(\mathsf{T}(\mathbf{M} \oplus \mathbf{M}^{*})\) 的），它约化为 \(d_{0}\) （在 \(\mathbf{A}\) 上）、为 \(d\) （在 \(\mathbf{M}\) 上）、为 \(d^{*}\) （在 \(\mathbf{M}^{*}\) 上）。限制 \(d_{\mathbf{J}}^{\mathbf{I}}\) ——即 \(\mathbf{D}\) 在子 A-模 \(\mathsf{T}_{\mathbf{J}}^{\mathbf{I}}(\mathbf{M})\) （ \(\mathsf{T}(\mathbf{M} \oplus \mathbf{M}^{*})\) 中的； \(\S 5\) ，第 6 号）上的限制——是 \(\mathsf{T}_{\mathbf{J}}^{\mathbf{I}}(\mathbf{M})\) 的一个 K-自同态，使得 \((d_{0}, d_{\mathbf{J}}^{\mathbf{I}}, d_{\mathbf{J}}^{\mathbf{I}})\) 是 \((\mathbf{A}, \mathsf{T}_{\mathbf{J}}^{\mathbf{I}}(\mathbf{M}), \mathsf{T}_{\mathbf{J}}^{\mathbf{I}}(\mathbf{M}))\) 的一个 K-导子。此外，对于 \(i \in \mathbf{I}\) ， \(j \in \mathbf{J}\) ，若记 \(\mathbf{I}' = \mathbf{I} - \{i\}\) ， \(\mathbf{J}' = \mathbf{J} - \{j\}\) ，则立即可验证，对收缩 \(c_{j}^{t}\) （ \(\S 5\) ，第 6 号）

\[
c _ {j} ^ {t} (d _ {J} ^ {\mathrm{I}} (z)) = d _ {J ^ {\prime}} ^ {\mathrm{I} ^ {\prime}} (c _ {j} ^ {t} (z)) \quad \text { for   all } z \in \mathbf {T} _ {J} ^ {\mathrm{I}} (\mathbf {M}).
\]

\begin{enumerate}
\def\labelenumi{(\arabic{enumi})}
\setcounter{enumi}{3}
\tightlist
\item
  设 \(M_{i}\) ( \(1 \leqslant i \leqslant 3\) ) 是三个 A-模，并且对每个 i 假设 \((d_{0}, d_{i}, d_{i})\) 是 \((\mathbf{A}, \mathbf{M}_{i}, \mathbf{M}_{i})\) 的一个导子；再对 n = 1 应用命题 16，则可导出一个导子 \((d_{0}, d_{ij}, d_{ij})\) （ \((\mathbf{A}, \mathbf{H}_{ij}, \mathbf{H}_{ij})\) 的，其中 \(\mathrm{H}_{ij} = \mathrm{Hom}_{\mathbf{A}}(\mathbf{M}_{i}, \mathbf{M}_{j})\) ），对每个有序对 \((i, j)\) 都是如此。利用这个记号，对 \(u \in \mathrm{Hom}_{\mathbf{A}}(\mathbf{M}_{1}, \mathbf{M}_{2})\) 和 \(v \in \mathrm{Hom}_{\mathbf{A}}(\mathbf{M}_{2}, \mathbf{M}_{3})\) ，
\end{enumerate}

\[
d _ {1 3} (v \circ u) = (d _ {2 3} v) \circ u + v \circ (d _ {1 2} u)\tag{39}
\]

正如从定义出发立即验证的那样。

\subsection{导子的泛问题；非交换情形}

在§10的其余部分中，所有代数均假定为结合的且有单位的，并且所有代数同态均假定为有单位的。

设 A 是一个 K-代数；张量积 \(A \otimes_{K} A\) 典范地具有一个 (A, A)-双模结构，在此结构下

\[
x. (u \otimes v). y = (x u) \otimes (v y)\tag{40}
\]

对所有 \(x, y, u, v\) （在 A 中）成立（ § 4，第 3 号，例 2）。对应于 A 中乘法的 K-线性映射 \(m: \mathrm{A} \otimes_{\mathbb{K}} \mathrm{A} \to \mathrm{A}\) （从而使得 \(m(x \otimes y) = xy\) ）是一个 (A, A)-双模同态；因此它的核 I 是 \(\mathrm{A} \otimes_{\mathbb{K}} \mathrm{A}\) 的一个子双模。

引理 1. 映射 \(\delta_{\mathbf{A}}: x \mapsto x \otimes 1 - 1 \otimes x\) 是 A 到 I 的一个 K-导子，且 I 作为左 A-模由 \(\delta_{\mathbf{A}}\) 的像生成。

第一个断言由以下事实得出：

\[
(x y) \otimes 1 - 1 \otimes (x y) = (x \otimes 1 - 1 \otimes x). y + x. (y \otimes 1 - 1 \otimes y)
\]

依 (40)。另一方面，若元素 \(\sum_{i} x_{i} \otimes y_{i}\) （ \(x_{i}, y_{i}\) 为 A 中的元素）属于 I，则按定义有 \(\sum_{i} x_{i} y_{i} = 0\) ，因此

\[
\sum_ {i} \left(x _ {i} \otimes y _ {i}\right) = \sum_ {i} x _ {i} (1 \otimes y _ {i} - y _ {i} \otimes 1)
\]

依 (40)，这就完成了引理的证明。

命题 17. 导子 \(\delta_{\mathbf{A}}\) 具有如下的泛性质：对每个 (A, A)-双模 E 以及每个 K-导子 \(d: \mathbf{A} \to \mathbf{E}\) ，存在且仅存在一个 (A, A)-双模同态 \(f: \mathbf{I} \to \mathbf{E}\) ，使得 \(d = f \circ \delta_{\mathbf{A}}\) 。

首先注意，对每个 (A, A)-双模同态， \(f: I \to E, f \circ \delta_{A}\) 是一个导子（第 7 号，命题 9）。反之，设 \(d: A \to E\) 是一个 K-导子；那么我们先证明：若存在 (A, A)-双模同态 \(f: I \to E\) 使得 \(d = f \circ \delta_{A}, f\) 由这一条件唯一确定，因为 \(\delta_{A}\) 的定义给出

\[
f (x \otimes 1 - 1 \otimes x) = d x
\]

而我们的断言由以下事实得出： \(\delta_{\mathbf{A}}\) 的像已经作为左 A-模生成 I（引理 1）：因此必然有

\[
f \left(\sum_ {i} x _ {i} \otimes y _ {i}\right) = \sum_ {i} x _ {i}. f (1 \otimes y _ {i} - y _ {i} \otimes 1) = - \sum_ {i} x _ {i}. d y _ {i}.
\]

反之，由于映射 \((x,y)\mapsto -x.dy\) （从 \(\mathbf{A}\times \mathbf{A}\) 到 \(\mathbf{E}\) 的）是 K-双线性的，因此存在且仅存在一个 K-线性映射 \(g:\mathbf{A}\otimes_{\mathbb{K}}\mathbf{A}\rightarrow \mathbf{E}\) ，使得 \(g(x\otimes y) = -x.dy\) ；只需验证限制 \(f\) （即 \(g\) 在 \(\mathbf{I}\) 上的限制）对左、右 A-模结构是 A-线性的。第一个断言是显然的，因为 \((xx') .dy = x.(x'.dy)\) ；为证明第二个断言，注意若 \(\sum_{i}x_{i}\otimes y_{i}\in \mathbf{I}\) 且 \(x\in \mathbf{A}\) ，则

\[
\sum_ {i} x _ {i}. d (y _ {i} x) = \sum_ {i} x _ {i}. d y _ {i}. x + \sum_ {i} (x _ {i} y _ {i}). d x
\]

但依 I 的定义有 \(\sum_{i}x_{i}y_{i} = 0\) ，这就完成了证明。

这样我们就定义了一个典范的 K-模同构 \(f \mapsto f \circ \delta_{A}\)

\[
\operatorname{Hom} _ {(\mathbf {A}, \mathbf {A})} (\mathbf {I}, \mathbf {E}) \rightarrow \mathbf {D} _ {\mathbf {K}} (\mathbf {A}, \mathbf {E})
\]

左边是由 A 到 E 的 (A, A)-双模同态构成的 K-模。
